\documentclass[11pt]{article}

\usepackage[T1]{fontenc}
\usepackage[utf8]{inputenc}
\usepackage{lmodern}
\usepackage{microtype}
\usepackage{amsmath,amssymb}
\usepackage[margin=1.1in]{geometry}
\usepackage{booktabs}
\usepackage{tabularx}
\usepackage{array}
\usepackage[colorlinks=true,linkcolor=blue,urlcolor=blue,citecolor=blue]{hyperref}

\title{Bayesian Inference for Nonlinear Models with a Genuine Zero Lower
Bound\\[0.6ex]\large A survey of exact and approximate sampling methods,
with two case studies}
\author{}
\date{August 2026}

\setcounter{tocdepth}{2}

\begin{document}
\maketitle

\begin{abstract}
\noindent
A zero (or effective) lower bound on nominal interest rates changes the
geometry of a Bayesian inference problem in ways that the word
``discontinuous'' hides more than it reveals. A censoring rule
$i_t=\max\{\ell,i_t^\star\}$ is continuous with a kinked gradient; a
deterministic solution map that selects between branches can create a genuine
density jump; a stochastic regime yields a mixed discrete--continuous
posterior; and a model whose equilibrium is not unique has no posterior at
all until a selection law is specified. These four geometries call for
different samplers, and misclassifying one as another is the most common
design error in this literature. We survey the linear benchmark (OccBin and
piecewise Kalman filtering), particle-filter likelihoods and particle MCMC,
Hamiltonian Monte Carlo on kinked and piecewise-smooth targets, event-driven
(reflection/refraction), discontinuous, and mixed variants of HMC, and
pseudo-marginal and differentiable-particle-filter methods, deriving the key
constructions rather than merely citing them. Two case studies apply the
classification: a two-country shadow-rate term-structure model, whose
hard-bound variant turns out to be a kink target with a tractable particle
likelihood; and a small New Keynesian model with an occasionally binding
policy rate, where the genuine discontinuity lives in the solution
correspondence (nonexistence and multiplicity of verified paths) rather
than in the max operator.
\end{abstract}

\tableofcontents
\newpage

\section{Introduction}\label{sec:intro}

In a linearized New Keynesian model, a zero lower bound (ZLB), or effective
lower bound (ELB), can be handled by solving two linear systems: one with the
policy constraint slack and one with the constraint binding. OccBin and its
successors iterate over a candidate sequence of regimes, apply a Kalman filter
conditional on that sequence, and verify the inequalities afterwards
(Guerrieri and Iacoviello 2015; Giovannini, Pfeiffer, and Ratto 2021). This
approach is valuable and fast. Its reach, however, is set by the conditional
Gaussian closure, which enters twice: once in the state transition and again
in the likelihood recursion. What it computes is therefore the likelihood of
a piecewise-linear Gaussian approximation, not of the nonlinear model.

The question we pursue in this survey is therefore narrower and harder: how
can we sample a posterior that retains an actual ZLB regime change in a
nonlinear state-space model, while keeping the transition kernel exact for the
model actually being estimated and making every approximation visible?

Before anything else, the word ``discontinuous'' needs care, because it is
routinely used for at least three different mathematical situations. The map
\[
 i_t=\max\{\ell,i_t^\star\}
\]
is continuous; it merely has a kink at \(i_t^\star=\ell\). A model with an
explicit regime variable \(r_t\in\{n,b\}\), or a solver that selects between
distinct equilibrium branches, can instead produce a mixed
discrete-continuous target, or a genuine jump in the density. These are three
different posterior geometries, and the modern Hamiltonian Monte Carlo (HMC)
literature addresses them with three different tools:
reflection/refraction HMC, discontinuous HMC, and mixed HMC. Treating all of
them as one kind of nonsmoothness is the first design error a practitioner
can make, and much of this survey is devoted to keeping them apart.

The survey has four goals. First, it derives the linear benchmark, the
piecewise-linear solution and its conditional Kalman likelihood, so that
the reader can see precisely what the fast standard method computes and what
it leaves out. Second, it derives the HMC kernels appropriate for
piecewise-smooth and mixed-support targets. Third, it derives the nonlinear
state-space likelihood and its particle approximation, which is what an exact
method must ultimately integrate. Fourth, it turns the comparison into
practice through two case studies: a shadow-rate term-structure model
(Section~\ref{sec:shadow}) and a small New Keynesian model with an
occasionally binding policy rate (Section~\ref{sec:nk}). Throughout, results
stated by a paper are cited to that paper; deductions of our own are written
in the notation introduced below and labelled as such.

\section{A common model and notation}\label{sec:model}

To compare methods that were developed in different literatures, we first fix
a single notation broad enough to contain them all. Let \(\vartheta\) denote
structural parameters, \(x_t\in\mathbb R^{d_x}\) the latent state, \(r_t\)
the ZLB regime, \(\epsilon_t\) a structural shock, and \(y_t\) the
observation. A general model is
\begin{equation}\label{eq:1}
\begin{aligned}
 r_t &= \mathcal R(x_{t-1},\epsilon_t,\vartheta),\\
 x_t &= F_{r_t}(x_{t-1},\epsilon_t,\vartheta),\\
 y_t &\sim g_{r_t}(y_t\mid x_t,\vartheta).
\end{aligned}
\end{equation}

For a monetary lower bound, define a shadow or notional nominal rate
\(i_t^\star=\iota(x_{t-1},\epsilon_t,\vartheta)\) and a bound \(\ell\). The
regime rule is
\begin{equation}\label{eq:2}
 r_t=n \iff i_t^\star>\ell,\qquad
 r_t=b \iff i_t^\star\leq\ell,
\end{equation}
and the observed policy rate may be \(i_t=i_t^\star\) in regime \(n\) and
\(i_t=\ell\) in regime \(b\). An equivalent complementarity form introduces a
multiplier \(\lambda_t\geq0\):
\begin{equation}\label{eq:3}
 i_t-\ell\geq0,\qquad \lambda_t\geq0,\qquad
 (i_t-\ell)\lambda_t=0.
\end{equation}
Equation \eqref{eq:3} says which inequality is true, but it does not by
itself say how the rest of the equilibrium is solved. That distinction
matters for inference: the likelihood is defined by the complete transition
and observation law, not by the bound in isolation.

The posterior over parameters, states, and regimes is
\begin{equation}\label{eq:4}
 \pi(\vartheta,x_{0:T},r_{1:T}\mid y_{1:T})
 \propto p(\vartheta)p(x_0\mid\vartheta)
 \prod_{t=1}^T p(r_t,x_t\mid x_{t-1},\vartheta)
 g_{r_t}(y_t\mid x_t,\vartheta).
\end{equation}

A measure-theoretic caution is required before \eqref{eq:4} can be used as a
density. The issue is that a density only exists relative to a dominating
measure, and a deterministic map need not have one with respect to Lebesgue
measure: if \(x_t\) is a deterministic function of \((x_{t-1},\epsilon_t)\)
and the shock has fewer dimensions than the state,
\(\dim\epsilon<\dim x\), then given \(x_{t-1}\) the state \(x_t\) is confined
to a lower-dimensional manifold, so its conditional law puts all its mass on
a Lebesgue-null set and no transition density exists; even when the
dimensions match, a change-of-variables Jacobian is required. The clean
primitive object is therefore not a density but the Markov transition kernel
\begin{equation}\label{eq:4a}
 K_\vartheta(x_{t-1},A,r)
 =\int \mathbf 1\{\mathcal R(x_{t-1},\epsilon,\vartheta)=r\}\,
       \mathbf 1\{F_r(x_{t-1},\epsilon,\vartheta)\in A\}\,
       \phi_\vartheta(d\epsilon),
\end{equation}
and \(p(r_t,x_t\mid x_{t-1},\vartheta)\) in \eqref{eq:4} denotes a density of
\(K_\vartheta\) with respect to an explicitly chosen dominating measure:
Lebesgue in \(x_t\) times counting in \(r_t\) when the shock map is a smooth
bijection onto its range. When it is not, the non-centered shock chart of
\eqref{eq:56} is the natural default coordinate system: there the target is a
density in \((\vartheta,x_0,\epsilon_{1:T})\) and no transition density is
needed at all. A similar bookkeeping point arises at the initial condition.
The factor \(p(x_0\mid\vartheta)\) in \eqref{eq:4} and the first-step factor
\(p(X_1\mid\vartheta)\) used by the particle recursion \eqref{eq:12a} must be
reconciled by adopting one of two conventions: either \(x_0\) is carried
explicitly in every particle and \eqref{eq:12a} begins from the prior
\(p(x_0\mid\vartheta)\) with a first transition draw, or the induced initial
law \(p(x_1\mid\vartheta)=\int K_\vartheta(x_0,\cdot)\,
p(x_0\mid\vartheta)\,dx_0\) is used and \(x_0\) is integrated out. Both are
correct; mixing them double-counts or drops the initial condition.

If the regimes are integrated out, the marginal likelihood is
\begin{equation}\label{eq:5}
 p(y_{1:T}\mid\vartheta)=
 \int p(x_0\mid\vartheta)
 \prod_{t=1}^T\sum_{r_t}p(r_t,x_t\mid x_{t-1},\vartheta)
 g_{r_t}(y_t\mid x_t,\vartheta)\,dx_{0:T}.
\end{equation}
The summation in \eqref{eq:5}, or the corresponding nonlinear integral, is
exactly the part that the piecewise Kalman filter avoids by conditioning on a
verified regime path. Keeping that summation in view is what distinguishes
the exact nonlinear methods of later sections from the linear benchmark of
Section~\ref{sec:linear}.

\subsection{Four different posterior geometries}\label{sec:geometries}

What determines the sampling problem is the measure in \eqref{eq:4}, not the
economic word ``regime''. Four cases must be kept separate, because each one
calls for a different sampler, and one of them calls for no sampler at
all, but for more modeling.

First, a censored rule such as \(i_t=\max\{\ell,i_t^\star\}\) is continuous.
If all other model equations and densities are continuous, the posterior is
also continuous, although its gradient may jump at the kink. Ordinary
Metropolized HMC remains a valid MCMC construction when the
nondifferentiable set has Lebesgue measure zero, but an integrator that
ignores the kink can be inefficient (Afshar and Domke 2015, p.~1). There is
no potential-energy jump to refract across.

Second, suppose a deterministic solution map produces different one-sided
densities on the two sides of a boundary. Write all continuous unknowns as
\(q\), let \(R(q)\) be the implied regime path, and let \(\pi_r(q)\) be the
density obtained from branch \(r\), where every \(\pi_r\) is expressed in one
named branch coordinate system; if different branches are naturally
parameterized in different coordinates, the change-of-variables Jacobian
into the common chart belongs inside \(\pi_r\). The joint density with
respect to Lebesgue measure in \(q\) and counting measure in \(r\) is
\begin{equation}\label{eq:5a}
 \Pi(q,r)=\pi_r(q)\,\mathbf 1\{r=R(q)\}.
\end{equation}
The indicator in \eqref{eq:5a} has a sharp consequence: at fixed \(q\), every
proposal \(r'\ne R(q)\) has zero target density. A regime-only Metropolis,
mixed-HMC, or Gibbs move therefore cannot cross this deterministic boundary.
One must instead move \(q\) across the boundary with an event-aware
continuous kernel, make a joint \((q,r)\) proposal whose density and Jacobian
are included, or integrate the states with a particle method. The one-sided
potential jump that event-driven dynamics must account for is
\begin{equation}\label{eq:5b}
 \Delta U(q_b)=-\log \pi_{r^+}(q_b^+)+\log \pi_{r^-}(q_b^-),
\end{equation}
where \(q_b^-\) and \(q_b^+\) denote limits from the two regions. For a pure
continuous \(\max\), these limits agree and \(\Delta U=0\), which recovers
the first case.

Third, if the statistical model assigns positive probabilities to both
regimes, say \(p(r_t\mid r_{t-1},x_{t-1},\vartheta)>0\) for both
values, then the posterior genuinely has mixed support. Mixed HMC and
particle Gibbs are valid candidates in this case, because a discrete move can
have positive target probability.

Fourth, if a nonlinear equilibrium solver can return several solutions, an
algorithmic tie-break such as ``first root found'' depends on solver order,
so it assigns the competing solutions no probabilities, and the target
remains undefined until the model supplies a solution-selection
distribution. No sampler can draw from a distribution that has not been
defined. We return to this issue at length in the New Keynesian
case study of Section~\ref{sec:nk}.

\section{What OccBin and piecewise Kalman filtering actually do}\label{sec:linear}

Before developing exact nonlinear methods, we set out the linear benchmark in
detail. Understanding precisely what the standard piecewise-linear machinery
computes, and where its two Gaussian closures enter, is what allows us
to say later, without hand-waving, what a nonlinear method adds.

\subsection{The two-regime linear system}\label{sec:linear-tworegime}

Guerrieri and Iacoviello (2015) write an occasionally binding constraint as a
reference regime \(n\) and an alternative regime \(b\). After linearization
around a common steady state, each regime has a linear rational-expectations
system. In reduced state-space notation, a candidate regime path \(r_{1:T}\)
gives
\begin{equation}\label{eq:6}
 x_t=T_{r_t}x_{t-1}+C_{r_t}+R_{r_t}\epsilon_t,
 \qquad \epsilon_t\sim N(0,Q),
\end{equation}
\begin{equation}\label{eq:7}
 y_t=Z_{r_t}x_t+D_{r_t}+u_t,
 \qquad u_t\sim N(0,H_{r_t}).
\end{equation}
The matrices come from solving each regime's equations as a linear
rational-expectations system, and the
candidate path is accepted only when the implied inequalities agree with the
regime labels; splicing arbitrary policy rules together would skip the
expectations consistency that the solution step enforces.

For a finite horizon, OccBin starts from a guessed regime path. It solves
backward from a terminal date at which the reference regime is imposed
forever, then simulates forward, checks the binding and slack inequalities,
and updates the path. Repeating this map gives a fixed point when it
converges. The method can have multiple fixed points; convergence of one
iteration is not a proof of global uniqueness (Guerrieri and Iacoviello 2015,
pp.~2--7).

\subsection{The conditional Kalman likelihood}\label{sec:linear-kalman}

Conditional on a verified regime path, the model \eqref{eq:6}--\eqref{eq:7}
is time-varying but Gaussian, so the likelihood is available in closed form
from the Kalman recursions. With filtered mean \(a_{t-1}\) and covariance
\(P_{t-1}\), prediction is
\begin{equation}\label{eq:8}
 a_t^- =T_{r_t}a_{t-1}+C_{r_t},
\qquad
 P_t^- =T_{r_t}P_{t-1}T_{r_t}^{\mathsf T}+R_{r_t}QR_{r_t}^{\mathsf T}.
\end{equation}
The innovation and its covariance are
\begin{equation}\label{eq:9}
 v_t=y_t-Z_{r_t}a_t^- -D_{r_t},
 \qquad
 F_t=Z_{r_t}P_t^-Z_{r_t}^{\mathsf T}+H_{r_t}.
\end{equation}
The update is
\begin{equation}\label{eq:10}
 K_t=P_t^-Z_{r_t}^{\mathsf T}F_t^{-1},\qquad
 a_t=a_t^-+K_tv_t,\qquad
 P_t=P_t^- -K_tF_tK_t^{\mathsf T}.
\end{equation}
For observed components \(\mathcal O_t\), the conditional log likelihood is
\begin{equation}\label{eq:11}
 \log p(y_{1:T}\mid r_{1:T},\vartheta)
 =-\frac12\sum_{t=1}^T
 \left[|\mathcal O_t|\log(2\pi)+\log|F_t|+v_t^{\mathsf T}F_t^{-1}v_t\right].
\end{equation}

Giovannini, Pfeiffer, and Ratto (2021, pp.~5--11) make the regime path
endogenous to the filter. Their piecewise Kalman filter (PKF) performs
\eqref{eq:8}--\eqref{eq:10}, estimates shocks and smoothed states, calls an
OccBin simulation to infer a new regime path, replaces the transition
matrices accordingly, and iterates. Failed fixed-point convergence is treated
as a bad parameter proposal. The publicly available Dynare implementation
follows this structure: a standard Kalman update is followed by an OccBin
solve, a new regime history, and a loop until the proposed and implied paths
agree. We note this as corroborating implementation detail, not as a
substitute for the paper's derivation.

\subsection{Exactness conditional on a linearization}\label{sec:linear-limits}

Exactness here is conditional in a strong sense. The PKF is exact for its
conditional time-varying linear-Gaussian model, but the model itself is a
piecewise-linear approximation. It omits precautionary
effects and higher-order nonlinearities that alter the probability of future
binding (Giovannini, Pfeiffer, and Ratto 2021, pp.~3--4, 26~ff.). OccBin also
needs a reference regime, a finite path iteration, and a rule for resolving
multiple fixed points. These are model and solver choices, not mere numerical
details.

The right role for the PKF in a nonlinear estimation exercise is that of a
\emph{limiting benchmark}: when the nonlinear model is replaced by its
verified piecewise-linear local representation, a correct nonlinear
implementation should reproduce \eqref{eq:11}, the regime path, and the
OccBin inequality checks. Passing this tie-out rules out a large class of
bookkeeping errors shared by the two implementations; validating the
nonlinear sampler on the model actually being estimated is the separate
task of the sections that follow.

\section{A nonlinear likelihood without regime smoothing}
\label{sec:particle}

What does it take to evaluate the marginal likelihood \eqref{eq:5} of the
genuinely nonlinear model \eqref{eq:1} without smoothing the regime away?
We answer in three stages: first with a generic particle filter, then with the
conditionally optimal particle filter built for piecewise-linear ZLB dynamics,
and finally with a family of branch-aware sigma-point constructions that sit
between a single Gaussian filter and a fully particle-based one.

\subsection{Generic particle filtering}
\label{sec:pf-generic}

Particle filtering approximates the sequence of filtering distributions by a
weighted swarm of samples, and as a by-product delivers an unbiased estimator
of the likelihood, the property on which all of the particle MCMC methods
of Section~\ref{sec:pmcmc} rest. For the nonlinear transition in
\eqref{eq:1}, we include the regime in the particle state when it is
stochastic; when it is deterministic, we compute it from each proposed state
and shock. Let $q_1(x_1\mid y_1,\vartheta)$ and
$q_t(x_t\mid x_{t-1},y_t,\vartheta)$ be proposal densities. At $t=1$, draw
$X_1^j\sim q_1$ and calculate
\begin{equation}\label{eq:12a}
 v_1^j=\frac{p(X_1^j\mid\vartheta)\,g(y_1\mid X_1^j,\vartheta)}
 {q_1(X_1^j\mid y_1,\vartheta)},\qquad
 W_1^j=\frac{v_1^j}{\sum_k v_1^k}.
\end{equation}
For $t\geq2$, resample an ancestor index
$A_t^j\sim\operatorname{Categorical}(W_{t-1}^{1:M})$, draw
$X_t^j\sim q_t(\cdot\mid X_{t-1}^{A_t^j},y_t,\vartheta)$, and set
\begin{equation}\label{eq:12b}
 v_t^j=\frac{p(X_t^j\mid X_{t-1}^{A_t^j},\vartheta)\,
                  g(y_t\mid X_t^j,\vartheta)}
 {q_t(X_t^j\mid X_{t-1}^{A_t^j},y_t,\vartheta)},\qquad
 W_t^j=\frac{v_t^j}{\sum_kv_t^k}.
\end{equation}
The resample-propagate likelihood estimator is
\begin{equation}\label{eq:13}
 \widehat p(y_{1:T}\mid\vartheta)=
 \prod_{t=1}^T\left(\frac1M\sum_{j=1}^M v_t^j\right).
\end{equation}

Why is \eqref{eq:13} unbiased? The argument is a single conditional
expectation, iterated backward. Fix the particle system at time $t-1$, that
is, condition on the particles $X_{t-1}^{1:M}$ and normalized weights
$W_{t-1}^{1:M}$. Each incremental weight $v_t^j$ depends on two fresh random
draws: the ancestor $A_t^j$ and the proposal draw $X_t^j$. Averaging first
over the proposal draw cancels the proposal density,
\[
 \mathbb E\!\left[v_t^j \,\middle|\, A_t^j=k\right]
 =\int \frac{p(x\mid X_{t-1}^{k},\vartheta)\,g(y_t\mid x,\vartheta)}
            {q_t(x\mid X_{t-1}^{k},y_t,\vartheta)}\,
       q_t(x\mid X_{t-1}^{k},y_t,\vartheta)\,dx
 =p(y_t\mid X_{t-1}^{k},\vartheta),
\]
the one-step predictive density from particle $k$. Averaging next over the
categorical ancestor draw gives
$\mathbb E[v_t^j]=\sum_k W_{t-1}^k\,p(y_t\mid X_{t-1}^{k},\vartheta)$, which
is exactly the particle approximation to the predictive integral
$p(y_t\mid y_{1:t-1},\vartheta)$ weighted by the time-$(t-1)$ filter.
Multiplying the per-period averages and repeating the conditional expectation
backward to $t=1$ telescopes the product into the Feynman--Kac normalizing
constant $p(y_{1:T}\mid\vartheta)$. Thus \eqref{eq:13} is nonnegative and
unbiased under the usual support and integrability conditions (Andrieu,
Doucet, and Holenstein 2010, Sec.~2.2; Aruoba et al. 2021, Sec.~6). One
caution: unbiasedness does not survive the logarithm. By Jensen's inequality
$\mathbb E[\log\widehat p]\leq\log p$, so the log-likelihood estimate is
biased downward, and this bias is what particle MCMC constructions must
handle at the level of the acceptance ratio rather than by averaging logs.

\subsection{The piecewise-linear conditionally optimal particle filter}
\label{sec:copf}

A generic bootstrap-style filter can be arbitrarily inefficient near the
bound, because blind proposals rarely land in the regime the data favor. The
conditionally optimal particle filter (COPF) of Aruoba et al. (2021) fixes
this for piecewise-linear ZLB dynamics by exploiting the regime structure in
the proposal itself. Their canonical form uses continuous piecewise-linear
decision rules,
\begin{equation}\label{eq:14}
 s_t=
 \begin{cases}
 \Phi_0(n)+\Phi_1(n)s_{t-1}+\Phi_\eta(n)\eta_t,&
 \eta_{1,t}<\zeta(s_{t-1}),\\[2pt]
 \Phi_0(b)+\Phi_1(b)s_{t-1}+\Phi_\eta(b)\eta_t,&
 \eta_{1,t}\geq\zeta(s_{t-1}),
 \end{cases}
\end{equation}
with a linear measurement equation. Conditional on a previous particle, the
observation density is then a mixture of two truncated Gaussian integrals.
Let $D_t^j(n)$ and $D_t^j(b)$ be the two regime contributions. Set
\begin{equation}\label{eq:15}
 \lambda_t^j=\frac{D_t^j(n)}{D_t^j(n)+D_t^j(b)}.
\end{equation}
The conditionally optimal proposal draws a regime with probability
$\lambda_t^j$, draws the structural innovation from the corresponding
truncated Gaussian, and assigns incremental weight
\begin{equation}\label{eq:16}
 \widetilde w_t^j=D_t^j(n)+D_t^j(b).
\end{equation}

The derivation is simply Bayes' rule. The optimal proposal is
\begin{equation}\label{eq:17}
 g_t^\star(s_t\mid y_t,s_{t-1})
 \propto p(y_t\mid s_t)\,p(s_t\mid s_{t-1}),
\end{equation}
so the importance ratio, which is the normalizing constant of \eqref{eq:17}, is
the predictive density $p(y_t\mid s_{t-1})$, which is exactly the sum in
\eqref{eq:16}. With vanishing measurement error, one, two, or no regime
solutions may explain a given observation; if every particle falls in the
no-solution case, the likelihood estimate is zero (Aruoba et al. 2021,
pp.~26--31).

This construction is the closest existing bridge to nonlinear ZLB inference:
it preserves an endogenous regime and avoids replacing the bound by a
softplus. Its limitations are equally important. The transition is still a
piecewise-linear approximation; the filter returns a likelihood value but no
parameter gradients, so an HMC kernel needs a separate gradient
construction; and when vanishing measurement error lets two regime
solutions explain the same observation, both contributions belong in the
mixture \eqref{eq:16}, since dropping one changes the predictive density
and hence the likelihood.

\subsection{A piecewise, truncated, and mixture UKF construction}
\label{sec:ukf}

Between a single Gaussian filter and a fully particle-based filter, the
unscented Kalman filter (UKF) literature contains a useful intermediate
family. Gaussian-sum filtering, which carries a weighted sum of Gaussian
components and updates each component with a Kalman-type step, goes back to
Alspach and Sorenson (1972). Within the sigma-point family, the truncated UKF
conditions the Gaussian approximation on bounded-support measurement
information (Garcia-Fernandez, Morelande, and Grajal 2012a), its mixture
variant applies the truncated update componentwise while retaining a Gaussian
mixture (Garcia-Fernandez, Morelande, and Grajal 2012b), and a UKF-aided
Gaussian-sum filter combines both ingredients (Gokce and Kuzuoglu 2015). We
cite these papers, and Alspach and Sorenson (1972), for orientation: we
were unable to obtain their full texts, so every equation in this section is
derived in our own notation rather than imported from them. Two
constrained-UKF papers we did study in full round out the picture: Kandepu,
Imsland, and Foss (2008, Secs.~II--III) project transformed sigma points onto
the feasible region, and Teixeira et al. (2010, Secs.~4--5, Table~1)
formulate interval-constrained unscented filtering, compare eight
projection-, constraint-, and truncation-based UKF variants, and note that
their motivating examples have multimodal densities. Together these methods
suggest a principled construction: preserve the binding and nonbinding
pieces explicitly, and defer exactness to a particle correction. None of
them alone evaluates the exact nonlinear-ZLB likelihood, which is why the
particle correction remains the final step.

\subsubsection{Why an ordinary UKF loses a branch}
\label{sec:ukf-branchloss}

To see what goes wrong with a single-Gaussian filter at the bound, let $z_t$
collect the uncertain variables that determine the next regime, for
example $z_t=(x_{t-1}^{\mathsf T},\epsilon_t^{\mathsf T})^{\mathsf T}$,
with a predictive approximation
\begin{equation}\label{eq:18}
 z_t\mid y_{1:t-1}\approx N(m_t,P_t).
\end{equation}
Let $D_b$ and $D_n$ denote the binding and nonbinding regions,
\begin{equation}\label{eq:19}
 D_b=\{z:b_t(z)\leq0\},
 \qquad
 D_n=\{z:b_t(z)>0\},
 \qquad
 b_t(z)=i_t^\star(z)-\ell .
\end{equation}
The exact predictive state density is a sum of integrals over these regions:
\begin{equation}\label{eq:20}
 p(x_t\mid y_{1:t-1})
 =\sum_{r\in\{b,n\}}
 \int_{D_r}p(x_t\mid z_t,r,\vartheta)\,
       p(z_t\mid y_{1:t-1})\,dz_t .
\end{equation}
Even when $p(z_t\mid y_{1:t-1})$ is Gaussian, \eqref{eq:20} is generally not
one Gaussian: conditioning on $D_r$ truncates the input, and the two maps
$F_b$ and $F_n$ can produce different means and covariances, so the output is
a two-component mixture at best. A conventional UKF instead propagates a
single sigma-point cloud through a branch function and replaces the
transformed distribution by its first two moments. That moment matching is a
legitimate approximation in smooth problems, but here it discards branch
mass, skewness, and multimodality exactly where the ZLB is informative. Nor
is the issue fixed by evaluating the $\max$ function at each sigma point:
the collapse to two moments happens regardless of how the points are mapped.

\subsubsection{Exact branch probabilities and truncated moments for an
affine boundary}
\label{sec:ukf-truncmoments}

Everything later in this section leans on one fully analytic building
block, the affine case, which we derive carefully.
Suppose
\begin{equation}\label{eq:21}
 b_t(z)=a^{\mathsf T}z-c,
 \qquad z\sim N(m,P),
 \qquad \sigma_b^2=a^{\mathsf T}Pa>0 .
\end{equation}
The scalar $u=a^{\mathsf T}z$ is $N(\mu_b,\sigma_b^2)$ with
$\mu_b=a^{\mathsf T}m$, so the branch event depends on $z$ only through one
Gaussian scalar. Define
\begin{equation}\label{eq:22}
 \gamma=\frac{c-\mu_b}{\sigma_b},
 \qquad
 \alpha_b=\Pr(b_t(z)\leq0)=\Phi(\gamma),
 \qquad
 \alpha_n=\Pr(b_t(z)>0)=1-\Phi(\gamma),
\end{equation}
where $\Phi$ and $\phi$ are the standard normal cdf and density.

To derive the conditional moments, standardize $v=(u-\mu_b)/\sigma_b$, so the
binding event is $v\leq\gamma$. The key elementary fact is that
$\frac{d}{dv}\bigl[-\phi(v)\bigr]=v\,\phi(v)$, so $\phi$ is its own
antiderivative up to the factor $v$. Direct integration then gives
\begin{equation}\label{eq:23}
 \mathbb E[v\mid v\leq\gamma]
 =\frac{\int_{-\infty}^{\gamma}v\,\phi(v)\,dv}{\Phi(\gamma)}
 =\frac{\bigl[-\phi(v)\bigr]_{-\infty}^{\gamma}}{\Phi(\gamma)}
 =-\frac{\phi(\gamma)}{\Phi(\gamma)}.
\end{equation}
For the second moment, integrate by parts, writing
$v^2\phi(v)=v\cdot v\phi(v)$ with antiderivative $-\phi$ for the second
factor:
\[
 \int_{-\infty}^{\gamma}v^2\phi(v)\,dv
 =\bigl[-v\,\phi(v)\bigr]_{-\infty}^{\gamma}
  +\int_{-\infty}^{\gamma}\phi(v)\,dv
 =\Phi(\gamma)-\gamma\,\phi(\gamma),
\]
so $\mathbb E[v^2\mid v\leq\gamma]=1-\gamma\lambda_b$ with
$\lambda_b=\phi(\gamma)/\Phi(\gamma)$, and subtracting the squared mean
$\lambda_b^2$ from \eqref{eq:23} gives
\begin{equation}\label{eq:24}
 \operatorname{Var}(v\mid v\leq\gamma)
 =1-\lambda_b(\lambda_b+\gamma),
 \qquad
 \lambda_b=\frac{\phi(\gamma)}{\Phi(\gamma)}.
\end{equation}
For the nonbinding event $v>\gamma$, the same two calculations over
$(\gamma,\infty)$ give
\begin{equation}\label{eq:25}
 \mathbb E[v\mid v>\gamma]=\lambda_n,
 \qquad
 \operatorname{Var}(v\mid v>\gamma)
 =1-\lambda_n(\lambda_n-\gamma),
 \qquad
 \lambda_n=\frac{\phi(\gamma)}{1-\Phi(\gamma)}.
\end{equation}

It remains to lift these scalar moments back to the full vector $z$. Since
$\operatorname{Cov}(z,u)=Pa$, the linear regression of $z$ on $u$ has
coefficient $Pa/\sigma_b^2$, and the Gaussian conditioning identity
decomposes $z$ into the scalar direction and an independent residual:
\begin{equation}\label{eq:26}
 z=m+\frac{Pa}{\sigma_b^2}(u-\mu_b)+\xi,
 \qquad
 \xi\perp u,
 \qquad
 \operatorname{Var}(\xi)=P-\frac{Paa^{\mathsf T}P}{\sigma_b^2}.
\end{equation}
Truncating the event $\{v\leq\gamma\}$ or $\{v>\gamma\}$ affects only the
$u$-term in \eqref{eq:26}; the residual $\xi$ is independent of the event and
keeps its moments. Substituting \eqref{eq:23}--\eqref{eq:25} into
\eqref{eq:26} therefore yields the exact first two moments of each truncated
branch:
\begin{equation}\label{eq:27}
 m_b=m-\frac{Pa}{\sigma_b}\lambda_b,
 \qquad
 P_b=P-\frac{Paa^{\mathsf T}P}{\sigma_b^2}\,
             \lambda_b(\lambda_b+\gamma),
\end{equation}
\begin{equation}\label{eq:28}
 m_n=m+\frac{Pa}{\sigma_b}\lambda_n,
 \qquad
 P_n=P-\frac{Paa^{\mathsf T}P}{\sigma_b^2}\,
             \lambda_n(\lambda_n-\gamma).
\end{equation}
We emphasize that \eqref{eq:22}, \eqref{eq:27}, and \eqref{eq:28} are the
exact moments of the original Gaussian law conditional on the branch event,
not a heuristic projection, and they are therefore the correct starting
point for a branch-aware sigma-point filter whenever the switching statistic
is affine in the augmented state and shock.

A word of practical advice before moving on: the Mills ratios
$\lambda_b=\phi(\gamma)/\Phi(\gamma)$ and
$\lambda_n=\phi(\gamma)/(1-\Phi(\gamma))$ underflow catastrophically in the
tails when evaluated from $\Phi$ directly, since numerator and denominator
both vanish at very different rates. The reader implementing these formulas
should work in log-cdf, log-survival, and log-density arithmetic (or use an
independently tested inverse-Mills-ratio routine), and should treat any
variance clamp added near $\gamma\to\pm\infty$ as a roundoff safeguard, not
as a substitute for stable probability evaluation.

For a nonlinear switching statistic $b_t$, no such closed form exists. One
can linearize $b_t$ at $m$ and reuse \eqref{eq:21}--\eqref{eq:28} with
$a=\nabla b_t(m)$ and $c=a^{\mathsf T}m-b_t(m)$, but that is an approximation
whose error grows with the curvature of the boundary inside the predictive
ellipsoid. It also inherits a sigma-point blind spot: a standard
$2d+1$-point set concentrated on the predictive mean can entirely miss a
low-probability binding region, so the estimated $\alpha_b$ can be exactly
zero even when the truth is small but positive. There are two defensible
responses: split components along the estimated boundary normal, which
refines the approximation, or apply the particle correction of
Section~\ref{sec:ukf-proposal}, which restores exactness. Either way, the
approximate and the exact ingredients should be kept clearly distinguished.

\subsubsection{The pure censored policy-rate likelihood}
\label{sec:ukf-censored}

A single censored Gaussian observation already exhibits the mixture
structure, with no UKF approximation at all, so we pause to work it out
completely. Let the shadow rate be
\begin{equation}\label{eq:29}
 i_t^\star\sim N(\mu,\sigma^2),
 \qquad i_t=\max\{\ell,i_t^\star\},
 \qquad y_t=i_t+u_t,
 \qquad u_t\sim N(0,V).
\end{equation}
The bound binds with probability
\begin{equation}\label{eq:30}
 \alpha_b=\Phi\!\left(\frac{\ell-\mu}{\sigma}\right),
 \qquad
 \alpha_n=1-\alpha_b.
\end{equation}
The censoring map $\max\{\ell,\cdot\}$ sends the entire lower tail of
$i_t^\star$ to the single point $\ell$, so with no measurement error the
distribution of $i_t$ has an atom of mass $\alpha_b$ at $\ell$ and a
continuous upper tail:
\begin{equation}\label{eq:31}
 p(i_t)=\alpha_b\,\delta_\ell(i_t)
 +\frac{1}{\sigma}\phi\!\left(\frac{i_t-\mu}{\sigma}\right)
  \mathbf 1\{i_t>\ell\}.
\end{equation}

The atom is often hidden when a Gaussian observation error is present, but it
does not disappear; it is merely convolved with the noise. To derive the
observed density, note first that on the binding event $i_t=\ell$ exactly, so
the binding contribution is $\alpha_b\,N(y_t;\ell,V)$. For the nonbinding
contribution we must integrate $N(y_t;i,V)$ against the upper Gaussian tail
of \eqref{eq:31}. The standard product-of-Gaussians identity does the work:
the product of two Gaussian densities in $i$ is, after collecting the
quadratic and linear terms in the exponent and completing the square, another
Gaussian in $i$ times a factor free of $i$,
\begin{equation}\label{eq:32}
 N(y_t;i,V)\,N(i;\mu,\sigma^2)
 =N(y_t;\mu,V+\sigma^2)\,N(i;\widetilde\mu_t,\widetilde\sigma^2),
\end{equation}
where matching coefficients gives the posterior-style precision-weighted
quantities
\begin{equation}\label{eq:33}
 \widetilde\sigma^2=\frac{\sigma^2V}{\sigma^2+V},
 \qquad
 \widetilde\mu_t=\frac{V\mu+\sigma^2y_t}{\sigma^2+V}.
\end{equation}
The factor $N(y_t;\mu,V+\sigma^2)$ pulls out of the integral, and what
remains is the mass that the Gaussian $N(\widetilde\mu_t,\widetilde\sigma^2)$
assigns to $(\ell,\infty)$. Integrating \eqref{eq:32} over $i>\ell$ therefore
gives the complete predictive density
\begin{equation}\label{eq:34}
 p(y_t)=
 \alpha_b\,N(y_t;\ell,V)
 +N(y_t;\mu,V+\sigma^2)
  \left[1-\Phi\!\left(
   \frac{\ell-\widetilde\mu_t}{\widetilde\sigma}\right)\right].
\end{equation}
The posterior probability that a given observation came from the binding
branch follows from Bayes' rule:
\begin{equation}\label{eq:35}
 \Pr(b\mid y_t)=
 \frac{\alpha_b\,N(y_t;\ell,V)}{p(y_t)}.
\end{equation}

Equation \eqref{eq:34} is the one-period analogue of the two-regime
truncated-Gaussian mixture underlying the COPF of
Section~\ref{sec:copf}. It also shows why replacing the max by a softplus is
not an innocent numerical trick: the softplus removes the atom in
\eqref{eq:31}, hence changes \eqref{eq:34}, and hence changes the likelihood
itself, not merely its evaluation. In a DSGE model, the other observables and
the state transition replace the scalar Gaussian factors in \eqref{eq:34},
but the same branch decomposition remains.

\subsubsection{Branchwise UKF update}
\label{sec:ukf-branchwise}

A full filtering step that carries both branches forward can now be
assembled from the affine building block of
Section~\ref{sec:ukf-truncmoments}. Suppose the filtering approximation before the next
observation is a mixture
\begin{equation}\label{eq:36}
 p(z_t\mid y_{1:t-1})
 \approx\sum_{j=1}^{J_{t-1}}\omega_{t-1,j}\,
 N(z_t;m_{t-1,j},P_{t-1,j}),
 \qquad \sum_j\omega_{t-1,j}=1.
\end{equation}
For each component $j$, calculate \eqref{eq:22}, \eqref{eq:27}, and
\eqref{eq:28}, obtaining a branch component $(j,r)$ with prior mass
\begin{equation}\label{eq:37}
 \rho_{j,r}=\omega_{t-1,j}\,\alpha_{j,r},
 \qquad r\in\{b,n\}.
\end{equation}
Use the branch-conditioned moments $(m_{j,r},P_{j,r})$ to form sigma points
$\chi_{j,r}^{(k)}$ with mean and covariance weights $w_m^{(k)},w_c^{(k)}$,
and propagate them through the branch-specific transition and observation
maps:
\begin{equation}\label{eq:38}
 X_{j,r}^{(k)}=F_r(\chi_{j,r}^{(k)},\vartheta),
 \qquad
 Y_{j,r}^{(k)}=h_r(X_{j,r}^{(k)},\vartheta).
\end{equation}
If process and observation noises are included in the augmented sigma point,
their covariances are already accounted for; otherwise add $Q_r$ and $V_r$
below. The branch prediction is
\begin{equation}\label{eq:39}
 \bar x_{j,r}=\sum_k w_m^{(k)}X_{j,r}^{(k)},
 \qquad
 P^x_{j,r}=\sum_k w_c^{(k)}
 (X_{j,r}^{(k)}-\bar x_{j,r})(X_{j,r}^{(k)}-\bar x_{j,r})^{\mathsf T}+Q_r,
\end{equation}
\begin{equation}\label{eq:40}
 \bar y_{j,r}=\sum_k w_m^{(k)}Y_{j,r}^{(k)},
\end{equation}
\begin{equation}\label{eq:41}
 S_{j,r}=\sum_k w_c^{(k)}
 (Y_{j,r}^{(k)}-\bar y_{j,r})(Y_{j,r}^{(k)}-\bar y_{j,r})^{\mathsf T}+V_r,
\end{equation}
\begin{equation}\label{eq:42}
 C_{j,r}=\sum_k w_c^{(k)}
 (X_{j,r}^{(k)}-\bar x_{j,r})(Y_{j,r}^{(k)}-\bar y_{j,r})^{\mathsf T}.
\end{equation}
The Gaussian UKF predictive likelihood, branch gain, and updated component
are
\begin{equation}\label{eq:43}
 \widehat L_{j,r}=N(y_t;\bar y_{j,r},S_{j,r}),
 \qquad K_{j,r}=C_{j,r}S_{j,r}^{-1},
\end{equation}
\begin{equation}\label{eq:44}
 m^+_{j,r}=\bar x_{j,r}+K_{j,r}(y_t-\bar y_{j,r}),
 \qquad
 P^+_{j,r}=P^x_{j,r}-K_{j,r}S_{j,r}K_{j,r}^{\mathsf T}.
\end{equation}
Bayes' rule combines the prior branch mass, the branch predictive
likelihood, and the evidence normalizer:
\begin{equation}\label{eq:45}
 \widetilde\omega_{j,r}=\rho_{j,r}\widehat L_{j,r},
 \qquad
 Z_t=\sum_{j=1}^{J_{t-1}}\sum_{r\in\{b,n\}}\widetilde\omega_{j,r},
 \qquad
 \omega^+_{j,r}=\frac{\widetilde\omega_{j,r}}{Z_t}.
\end{equation}
Thus the branch-aware UKF approximation after observing $y_t$ is
\begin{equation}\label{eq:46}
 p(x_t\mid y_{1:t})
 \approx\sum_{j,r}\omega^+_{j,r}\,
 N(x_t;m^+_{j,r},P^+_{j,r}).
\end{equation}

If one insists on collapsing \eqref{eq:46} to a single Gaussian, the
moment-preserving collapse is
\begin{equation}\label{eq:47}
 \bar m_t=\sum_{j,r}\omega^+_{j,r}m^+_{j,r},
\end{equation}
\begin{equation}\label{eq:48}
 \bar P_t=\sum_{j,r}\omega^+_{j,r}
 \left[P^+_{j,r}+
 (m^+_{j,r}-\bar m_t)(m^+_{j,r}-\bar m_t)^{\mathsf T}\right].
\end{equation}
The second term in \eqref{eq:48} is the between-branch variance, and it does
real work: dropping it makes the next period's
threshold probability too small or too large whenever the binding and
nonbinding components are well separated. Keeping all components, on the
other hand, lets the mixture grow to as many as $2^t$ histories (the
classical Gaussian-sum growth problem), so practical filters use pruning,
merging, or an IMM-style moment collapse. Each of these operations changes
the filtering density, so its effect belongs in the error budget of the run
alongside the moment-matching step itself.

\subsubsection{Relation to IMM-UKF and to deterministic ZLB support}
\label{sec:ukf-imm}

It is natural to ask how the branchwise construction relates to the classical
multiple-model filters, and the answer depends critically on whether the
regime is genuinely stochastic. For a stochastic Markov regime, the preceding
construction becomes a multiple-model UKF. If
$p(r_t=r\mid r_{t-1}=s)=\Pi_{sr}$, the mode-prediction probability is
\begin{equation}\label{eq:49}
 c_{t,r}=\sum_s\Pi_{sr}\mu_{t-1,s},
 \qquad
 \mu_{t-1,s\mid r}=\frac{\Pi_{sr}\mu_{t-1,s}}{c_{t,r}},
\end{equation}
and the IMM mixing mean and covariance supplied to the $r$-specific UKF are
\begin{equation}\label{eq:50}
 m_{t-1\mid r}=\sum_s\mu_{t-1,s\mid r}m_{t-1,s},
\end{equation}
\begin{equation}\label{eq:51}
 P_{t-1\mid r}=\sum_s\mu_{t-1,s\mid r}
 \left[P_{t-1,s}+
 (m_{t-1,s}-m_{t-1\mid r})(m_{t-1,s}-m_{t-1\mid r})^{\mathsf T}\right].
\end{equation}
After the branchwise update, the mode probability is proportional to
$c_{t,r}\widehat L_{t,r}$. Equation \eqref{eq:49} is Bayes' rule on the mode
chain, and \eqref{eq:50}--\eqref{eq:51} apply the moment-preserving collapse
\eqref{eq:47}--\eqref{eq:48} within each destination mode. This is the
interacting-multiple-model (IMM) logic introduced for Markovian switching
systems by Blom and Bar-Shalom (1988), here with UKF components in place of
linear Kalman components. A constrained multiple-model UKF we examined
applies the same Markov-jump structure with truncated noise supports and
model-conditioned feasible regions (Zhang et al. 2020, Secs.~II--III).

The deterministic ZLB is different, and the difference is exactly the
support issue of \eqref{eq:5a}. There is no free $\Pi_{sr}$ to tune: the
branch is $r_t=R(z_t)$, and at a fixed $z_t$ the other label has zero
support. Equations \eqref{eq:49}--\eqref{eq:51} are therefore a valid
approximation only if we deliberately introduce a stochastic regime law. For
the original complementarity model, $\alpha_{j,r}$ in \eqref{eq:37} must be
computed from the continuous predictive law, and the regime label must remain
a derived quantity rather than a free state.

\subsubsection{UKF proposals inside a corrected particle filter}
\label{sec:ukf-proposal}

Every construction in Section~\ref{sec:ukf} up to this point approximates the
likelihood. The branchwise UKF becomes much more defensible when it proposes
particles but does not define the likelihood by itself: the proposal can be
crude, because importance weighting repairs it. Let
$q_{t,j,r}(z_t\mid y_t)$ be a branch-conditioned UKF proposal supported on
$D_r$, and let $\alpha_{j,r}$ be the branch-selection probability. Draw $r$
with probability $\alpha_{j,r}$, then draw $z_t\sim q_{t,j,r}$. For a
stochastic regime, the incremental importance weight is
\begin{equation}\label{eq:52}
 w_t=
 \frac{p(r_t\mid r_{t-1},x_{t-1},\vartheta)\,
       p(z_t\mid x_{t-1},r_t,\vartheta)\,
       g_{r_t}(y_t\mid x_t,\vartheta)}
      {\alpha_{j,r_t}\,
       q_{t,j,r_t}(z_t\mid x_{t-1},y_t,\vartheta)}.
\end{equation}
For a deterministic regime, $r_t=R(z_t)$ is not an independently sampled
variable. If the shock-augmented transition is used, the corresponding weight
is
\begin{equation}\label{eq:53}
 w_t=
 \frac{p(\epsilon_t\mid\vartheta)\,
       g_{R(z_t)}(y_t\mid x_t,\vartheta)}
      {\alpha_{j,R(z_t)}\,
       q_{t,j,R(z_t)}(z_t\mid x_{t-1},y_t,\vartheta)},
 \qquad x_t=F_{R(z_t)}(x_{t-1},\epsilon_t,\vartheta).
\end{equation}
Notice what does and does not enter these weights: the exact numerator and
the proposal density determine $w_t$, while the UKF's Gaussian moment
approximation appears only inside the proposal, where any error is corrected
by the weight.

One coordinate subtlety is implicit in \eqref{eq:53} and needs spelling out.
The numerator density is written in the shock coordinate $\epsilon_t$, while
the proposal density is written in the augmented coordinate $z_t$. The two
are consistent only if either (a) the proposal draws the shock directly, so
that $q$ is a density in $\epsilon_t$ (with $x_{t-1}$ fixed, the map
$\epsilon_t\mapsto z_t$ is an affine embedding and both densities live in the
same chart), or (b) the proposal density in $z_t$ is converted to the
$\epsilon_t$ chart with the absolute Jacobian of the transformation. We
recommend convention (a): propose shocks, derive states. Under either
convention, and under the usual support and integrability conditions,
inserting \eqref{eq:52} or \eqref{eq:53} into the particle estimator
\eqref{eq:13} preserves the nonnegative unbiased likelihood property of
Section~\ref{sec:pf-generic}. A UKF-only filter has no such correction; it
remains an approximate Gaussian-sum likelihood.

A precedent for this architecture is the iterative truncated unscented
particle filter of Wang et al. (2020, Secs.~3.1--3.2, Algorithm~1). There an
iterated-UKF update is truncated to the constraint region, the first two
moments of the truncated density define a Gaussian proposal, and each
particle is reweighted by the ratio of target factors to the proposal
density. Two differences from \eqref{eq:52}--\eqref{eq:53} matter for the
ZLB. First, their constraint region is exogenous, whereas the ZLB branch is
selected by the model itself, which is why the branch probability
$\alpha_{j,r}$ and the branch-specific maps $F_r,h_r$ appear in our weights.
Second, their weight numerator uses a Gaussian approximation of the
predictive density (their eq.~36), not the exact transition density, so
their filter is a proposal-quality benchmark rather than an exactly weighted
likelihood estimator; a filter whose likelihood estimate is to be trusted as
unbiased must keep the exact shock density in the numerator, as in
\eqref{eq:53}.

The hierarchy that emerges is worth stating plainly. A split or truncated
mixture UKF is a fast, useful diagnostic and a good proposal engine. A
branch-corrected particle filter is the authoritative nonlinear likelihood.
The COPF construction of Section~\ref{sec:copf} is the exact
piecewise-linear benchmark. None of this, however, turns the UKF proposal
into an exact source of HMC gradients: an HMC scheme must target either the
corrected extended particle state or an explicitly acknowledged
approximation (Sections~\ref{sec:hmc} and~\ref{sec:pmcmc}), and the
shadow-rate case study of Section~\ref{sec:shadow} shows what these
constructions become when the switching lives in the observation map
rather than in the transition.

\section{Ordinary HMC and how the lower bound changes its geometry}\label{sec:hmc}

Hamiltonian Monte Carlo is the natural starting point for posterior sampling
in high-dimensional latent-variable models, and understanding precisely where
an occasionally binding constraint breaks its assumptions is the organizing
question for the rest of this survey. For a smooth posterior
$\pi(q)\propto e^{-U(q)}$, introduce momentum $p\sim e^{-K(p)}$, form the
Hamiltonian $H=U+K$, and evolve the phase point by Hamilton's equations
(Neal 2011, Sec.~5.2)
\begin{equation}\label{eq:54}
 \dot q=\nabla_pK(p),\qquad \dot p=-\nabla_qU(q).
\end{equation}
With Gaussian momentum $K(p)=\frac12p^{\mathsf T}M^{-1}p$, the leapfrog
integrator discretizes \eqref{eq:54} as
\begin{equation}\label{eq:55}
 p\leftarrow p-\tfrac\epsilon2\nabla U(q),\quad
 q\leftarrow q+\epsilon M^{-1}p,\quad
 p\leftarrow p-\tfrac\epsilon2\nabla U(q).
\end{equation}
The leapfrog map is reversible and volume preserving, so a final Metropolis
correction with acceptance probability $1\wedge e^{-\Delta H}$ exactly
accounts for the integration error.

An attractive way to deploy HMC on a nonlinear state-space model is to
sample the parameters and the latent shocks jointly, sidestepping the filter
entirely. Childers et al.\ (2022, pp.~10--16) pursue this idea: for a state
equation $x_t=F(x_{t-1},\epsilon_t,\vartheta)$, the joint log posterior is,
up to a constant,
\begin{equation}\label{eq:56}
 \log p(\vartheta,x_0,\epsilon_{1:T}\mid y)
 =\log p(\vartheta)+\log p(x_0\mid\vartheta)
 +\sum_{t=1}^T\left[\log p(\epsilon_t\mid\vartheta)
 +\log g(y_t\mid x_t,\vartheta)\right].
\end{equation}
This can remove a costly nonlinear filter, but the price is that HMC must
differentiate the complete simulation map $F$ composed over $T$ periods. Here
the lower bound intrudes. A hard $\max$ is only piecewise differentiable:
automatic differentiation happily supplies the derivative of whichever branch
is active, but it reports no boundary event and no information about the
kink. A branch-dependent equilibrium solver may in addition jump or be
multi-valued across the boundary, and discrete resampling inside a particle
filter changes ancestor indices discontinuously. Replacing the bound with a
softplus restores smoothness, but by removing the censoring atom it changes
the likelihood itself (Section~\ref{sec:ukf-censored} computes this exactly),
so it defines a separate approximate model rather than a smooth route to the
constrained one.

\subsection{Validity of Metropolized HMC on a kinked target}\label{sec:kink-validity}

What does a kink, a continuous potential with a
discontinuous gradient, actually cost? In the two-country shadow-rate
case study of Section~\ref{sec:shadow}, the target turns out to be continuous
with a kinked gradient rather than discontinuous, so the validity of ordinary
leapfrog HMC on such targets deserves a careful statement instead of a
passing remark. We state and sketch the following, since we could not locate
a reference at this level of generality; the assumptions are chosen to cover
a continuous piecewise-quadratic-plus-smooth potential, which is exactly what
a piecewise-affine observation map inside a Gaussian likelihood produces.

\paragraph{Assumptions.}
$U$ is continuous on $\mathbb R^d$; there is a closed Lebesgue-null set $N$
arising from a finite (or locally finite) piecewise-regular partition of
$\mathbb R^d$ (for the target of Section~\ref{sec:shadow}, a finite union
of hyperplanes) such that $U$ is continuously differentiable on
$\mathbb R^d\setminus N$ with $\nabla U$ locally bounded, the branch gradient
on each piece is deterministic and single-valued, and a fixed boundary
convention assigns a value on $N$; and $e^{-U}$ is integrable. These
assumptions are \emph{not} satisfied merely because a potential is
almost-everywhere differentiable: a potential evaluated through an implicit
multi-root equilibrium solver, a state-dependent iteration count, or any
branch rule that is not a fixed partition of $\mathbb R^d$ is outside this
proposition until a separate theorem covers it.

\paragraph{Claim.}
Fix a step size $\epsilon$ and a step count $L$. The set of initial phase
points $(q,p)$ whose leapfrog trajectory \eqref{eq:55} evaluates $\nabla U$
at a point of $N$ has Lebesgue measure zero. Off this null set, the $L$-step
leapfrog map is well defined, volume preserving, and reversible in the usual
momentum-flip sense, so the Metropolized chain that accepts with probability
$1\wedge e^{-\Delta H}$ leaves $\pi\propto e^{-U}$ invariant.

\emph{Sketch.} Each half-kick $p\mapsto p-\tfrac\epsilon2\nabla U(q)$ at
fixed $q\notin N$ and each drift $q\mapsto q+\epsilon M^{-1}p$ is a
measure-preserving shear. The set of starting points that land on $N$ at a
given substep is the preimage of a null set under a composition of
measure-preserving maps defined off a null set, hence null; a finite union
over the $2L+1$ substeps is null. Reversibility and volume preservation on
the remaining full-measure set follow from the standard shear argument (Neal
2011, Sec.~5.2), which nowhere uses smoothness across $N$. The
piecewise-regular partition assumption is what makes the shear composition
well defined off a null set; the argument is \emph{not} valid for an
arbitrary nonsmooth composition merely because individual shears preserve
measure. $\square$

Validity is not efficiency. Crossing the kink inside a drift substep replaces
the usual $O(\epsilon^3)$ local energy error by
\begin{equation}\label{eq:56a}
 \Delta H \;=\; O\!\left(\epsilon\,
 \bigl\|\nabla U^{+}-\nabla U^{-}\bigr\|\right)
\end{equation}
at the crossing, where $\nabla U^{\pm}$ are the one-sided gradients, so
acceptance degrades with the frequency of kink crossings and the size of the
gradient jump. This is the inefficiency, not invalidity, observed by
Afshar and Domke (2015, p.~1). Two consequences follow. First, the reflection
and refraction machinery of Section~\ref{sec:event} degenerates on a kink:
the potential jump in \eqref{eq:58} is $\Delta U=0$, so \eqref{eq:59} leaves
the momentum unchanged, and event detection buys accuracy, never
correctness. Second, for a continuous \emph{piecewise-quadratic} potential,
the Hamiltonian flow is available in closed form on each piece; this is the
construction of exact HMC for truncated multivariate Gaussians (Pakman and
Paninski 2014, which we cite for orientation; we were unable to obtain the
full text), and it is directly relevant to sampling the cell-restricted
Gaussians of Section~\ref{sec:shadow-filter}. A boundary-exact integrator is
therefore an optimization for the kinked case and a requirement only for a
genuine density jump.

\section{The UKF filtering collapse problem and marginal-likelihood gradient discontinuities}\label{sec:ukf-gradient}

The preceding section established that ordinary leapfrog HMC remains valid on
continuous-kink targets, albeit with degraded efficiency. This section
addresses a more insidious geometry that arises when a Gaussian filter
(Kalman, extended Kalman, or unscented) is used to compute the marginal
likelihood $p(y_{1:T}\mid\vartheta)$ in a model with occasionally binding
constraints. The composed map
$\vartheta\mapsto\log p(y_{1:T}\mid\vartheta)$ can develop steep or
discontinuous gradients even when the underlying economic model uses a smooth
softplus constraint rather than a hard max, because the filter's variance
collapse near the bound amplifies parameter-space gradient jumps. This is the
geometry encountered in the two-country shadow-rate model of
Section~\ref{sec:shadow} when a UKF replaces the particle filter, and it is
the primary obstacle to making gradient-based MCMC tractable for realistic
DSGE models with zero lower bounds.

\subsection{Why UKF variance collapses at the bound}\label{sec:ukf-collapse-mechanism}

Consider a shadow nominal rate $i_t^\star$ that is a deterministic function of
a latent state $x_t$, with a smoothed constraint
\begin{equation}\label{eq:ukf1}
 i_t = \ell + \operatorname{softplus}(i_t^\star - \ell, \alpha)
     = \ell + \frac{1}{\alpha}\log(1 + e^{\alpha(i_t^\star-\ell)}),
\end{equation}
and suppose $i_t$ is observed with small Gaussian noise. When the data pin
$i_t$ near $\ell$ for many consecutive periods, a UKF (or any Gaussian filter)
updates the filtering distribution over $x_t$ by conditioning on those
near-bound observations. The Kalman gain and innovation covariance depend on
the predictive variance; when the observations are informative and consistent,
the filtering variance shrinks. For a state whose shadow rate is deeply
negative, all values below a threshold produce approximately the same observed
rate $i_t\approx\ell$, so the likelihood provides no information about
\emph{how far} below the bound the shadow rate is. The filtering distribution
thus develops a one-sided collapse: variance in the direction that affects the
observed rate shrinks to near zero, while variance in the orthogonal
``shadow-depth'' direction can grow or remain large, depending on the
transition dynamics.

The softplus in \eqref{eq:ukf1} does not prevent this collapse; it merely
makes the transition from the informative region ($i_t^\star\approx\ell$) to
the flat plateau ($i_t^\star\ll\ell$) continuous rather than abrupt. As the
softplus temperature $\alpha$ increases, the transition region becomes
steeper, and the geometry approaches that of a hard max.

\subsection{How variance collapse creates parameter-space gradient jumps}\label{sec:collapse-gradient}

A UKF propagates a deterministic set of sigma points through the nonlinear
transition and observation maps. When the filtering variance is large, the
sigma points are spread out; small parameter perturbations typically move all
points by small amounts, and the marginal likelihood gradient
$\nabla_\vartheta\log p(y\mid\vartheta)$ varies smoothly. When the variance
collapses, the sigma points cluster tightly near the filtering mean. A small
parameter change can now shift the entire cluster across the steep part of the
softplus, causing all sigma points to change regime contribution
simultaneously. The result is a sharp change in the predicted observations,
innovation covariance, and log-likelihood contribution, which manifests as a
large jump in $\nabla_\vartheta\log p(y\mid\vartheta)$.

This is a \emph{continuous-kink geometry}: the likelihood
$\log p(y\mid\vartheta)$ is continuous (by continuity of the softplus and the
Gaussian observation density), but its gradient can have jumps of size
$O(\alpha)$ when a collapsed sigma-point cloud crosses the steep softplus
region. The gradient jump is not a modeling artifact; it is a consequence of
the variance collapse in the filter combined with the deterministic sigma-point
rule.

\subsection{Why this differs from a value jump}\label{sec:kink-vs-jump-ukf}

The UKF filtering collapse creates a continuous-kink target, not a value-jump
target. At a parameter-space boundary where the sigma points cross the steep
softplus region, the one-sided likelihood values match:
\[
 \lim_{\vartheta\to\vartheta_b^-}\log p(y\mid\vartheta)
 = \lim_{\vartheta\to\vartheta_b^+}\log p(y\mid\vartheta),
\]
but the gradients do not:
\[
 \lim_{\vartheta\to\vartheta_b^-}\nabla\log p(y\mid\vartheta)
 \ne \lim_{\vartheta\to\vartheta_b^+}\nabla\log p(y\mid\vartheta).
\]

By the argument of Section~\ref{sec:kink-validity}, ordinary Metropolized HMC
remains \emph{valid} for this target: the leapfrog map is well defined almost
everywhere, reversible, and volume preserving, so the Metropolis correction
targets the correct posterior. However, leapfrog is \emph{inefficient}: each
trajectory that crosses a steep-gradient region incurs an $O(\epsilon\,\Delta
g)$ energy error, where $\Delta g$ is the gradient-jump magnitude and
$\epsilon$ is the step size. When $\Delta g$ is large (high softplus
temperature $\alpha$ or severe variance collapse), acceptance rates degrade,
and the chain can become stuck near boundaries or in one filtering regime.

This is the practical obstacle: the method is mathematically correct but
computationally poor. Reflection/refraction HMC (Section~\ref{sec:event}) does
not help, because there is no potential jump $\Delta U$ to pay with kinetic
energy—the target is continuous. What is needed instead is either a different
filter that maintains variance (particle methods), or an event-aware integrator
that detects steep-gradient crossings and adjusts the trajectory accordingly.

\subsection{Why integer regime solvers create value jumps}\label{sec:integer-regime-jumps}

The analysis above assumed a smoothed constraint. When the model instead uses
an integer solver that selects discrete regime paths—for example, a
quantitative-easing on/off indicator or an OccBin-style spell-duration
selector—the geometry changes fundamentally. The solver maps parameters to a
discrete regime path $r_{1:T}\in\{0,1\}^T$ by checking equilibrium consistency
inequalities. At a parameter-space boundary where the solver switches from path
$j$ to path $k$, the regime-dependent state-space matrices $(F_{r_t},
Q_{r_t}, H_{r_t})$ change discontinuously, even though the parameter
$\vartheta$ changes smoothly.

The filtered likelihood is then a composition:
\[
 \vartheta \xrightarrow{\text{solver}} r_{1:T}
          \xrightarrow{\text{state-space}} (F_t, Q_t, H_t)
          \xrightarrow{\text{UKF}} \log p(y\mid\vartheta).
\]
When the regime path jumps, the UKF evaluates a different piecewise-linear
system. The innovation sequence, covariance updates, and cumulative
log-likelihood all depend on the selected regime, so the one-sided likelihood
values generally do not match:
\[
 \lim_{\vartheta\to\vartheta_b^-}\log p(y\mid\vartheta, r^-)
 \ne \lim_{\vartheta\to\vartheta_b^+}\log p(y\mid\vartheta, r^+).
\]

This is a \emph{value-jump geometry}, the case for which reflection/refraction
HMC and discontinuous HMC were designed. The potential has a finite jump
$\Delta U = -[\log p(y\mid\vartheta_b,r^+) - \log p(y\mid\vartheta_b,r^-)]$,
and standard leapfrog is invalid: a trajectory that crosses the boundary
without an event map systematically under- or over-estimates the Hamiltonian,
leading to biased acceptance and wrong stationary distribution.

\subsection{Measurement protocol for UKF-based posteriors}\label{sec:ukf-measurement}

Before choosing an HMC extension for a UKF-filtered ZLB model, the actual
boundary geometry must be measured. The following protocol distinguishes
continuous kinks from value jumps and quantifies gradient variation.

\paragraph{Step 1: Identify candidate boundaries.}
Run a short HMC warm-up and record parameter values where:
\begin{itemize}
\item the softplus input $i_t^\star - \ell$ is near zero (within a few
      multiples of $1/\alpha$);
\item the UKF filtering variance for the shadow-rate state is small (orders of
      magnitude below the prior variance);
\item an integer regime solver (if present) changes its selected path.
\end{itemize}

\paragraph{Step 2: Value-continuity test.}
For parameter pairs $(\vartheta^-, \vartheta^+)$ straddling a candidate
boundary, evaluate:
\[
 \Delta\ell = |\log p(y\mid\vartheta^+) - \log p(y\mid\vartheta^-)|.
\]
If $\Delta\ell / |\log p(y\mid\vartheta^-)| < 10^{-6}$, classify as
continuous kink. If $\Delta\ell = O(1)$ or larger, classify as value jump.

\paragraph{Step 3: Gradient-jump magnitude.}
For continuous-kink boundaries, evaluate the one-sided gradients and compute:
\[
 \Delta g = \|\nabla_\vartheta\log p(y\mid\vartheta^+)
             - \nabla_\vartheta\log p(y\mid\vartheta^-)\|_2.
\]
Large $\Delta g$ relative to the typical gradient norm indicates poor leapfrog
geometry.

\paragraph{Step 4: Boundary crossing rate.}
Run HMC trajectories of length $L$ steps and record:
\begin{itemize}
\item fraction of trajectories that cross a steep-gradient or regime-switch
      boundary;
\item acceptance rate conditional on crossing vs not crossing;
\item average Hamiltonian error on crossing vs non-crossing trajectories.
\end{itemize}
High crossing rate combined with low crossing acceptance indicates the boundary
is a bottleneck.

\paragraph{Step 5: Gradient Lipschitz estimation.}
Sample parameter pairs with $\|\vartheta_1 - \vartheta_2\| = \delta$ and
estimate:
\[
 L_{\text{est}} = \max_{\text{pairs}}
   \frac{\|\nabla\log p(y\mid\vartheta_1) - \nabla\log p(y\mid\vartheta_2)\|}
        {\|\vartheta_1 - \vartheta_2\|}.
\]
For a softplus model, $L_{\text{est}}$ grows approximately linearly with the
temperature $\alpha$.

\subsection{Why particle filters smooth the gradient}\label{sec:pf-smoothing}

A particle filter approximates the filtering distribution by a weighted cloud
of samples. Unlike a UKF's deterministic sigma points, particles retain
diversity even when the likelihood strongly favors one region. When some
particles have shadow rates slightly above the bound and others slightly below,
the marginal likelihood is a mixture:
\[
 p(y_t\mid y_{1:t-1},\vartheta)
 \approx \sum_{i=1}^M W_{t-1}^i\, p(y_t\mid X_{t-1}^i,\vartheta),
\]
where each particle $X_{t-1}^i$ contributes its own predictive density. A
small parameter change moves each particle slightly, changing the mixture
weights smoothly rather than flipping an entire cluster. The gradient
$\nabla_\vartheta\log p(y\mid\vartheta)$ is then an average over particle
contributions, and the averaging smooths out the sharp transitions that a
collapsed UKF exhibits.

This is the fundamental reason differentiable particle filters (LEDH-PFPF-OT
and similar constructions) are pursued for ZLB models: they maintain diversity
by design, which bounds the gradient variation and makes HMC geometrically
tractable. The computational cost is the price of that smoothness.

\subsection{Candidate solutions for the UKF gradient problem}\label{sec:ukf-solutions}

Given the measurement protocol of Section~\ref{sec:ukf-measurement}, the
following strategies apply depending on the diagnosed geometry.

\paragraph{For continuous-kink targets (softplus ZLB, no integer regime).}
\begin{enumerate}
\item \textbf{Smaller step size.} Reduce $\epsilon$ to lower the per-crossing
      energy error \eqref{eq:56a}. This is the simplest fix but increases
      cost per effective sample.
\item \textbf{Event-aware integration (Tran and Kleppe 2025).} Detect when a
      leapfrog step would cross a steep-gradient region (by monitoring the
      softplus input or a related margin), truncate the step at the boundary,
      recompute the gradient on the new side, and continue. This is
      computationally intensive (requires bisection and re-evaluation) but
      reduces integration error.
\item \textbf{Neural surrogate proposal.} Train a neural network
      $s_\phi(\vartheta)$ to approximate $\nabla_\vartheta\log p(y\mid\vartheta)$
      using automatic differentiation through the UKF at training points. Use
      $s_\phi$ to drive leapfrog proposals, but evaluate the \emph{true} UKF
      likelihood for Metropolis acceptance. A poor surrogate leads to low
      acceptance, not wrong posterior (Section~\ref{sec:pmcmc}).
\item \textbf{UKF variance inflation.} Artificially inflate the predicted
      covariance $P_{t|t-1}$ when the variance is below a threshold and the
      state is near a bound. This prevents full collapse and smooths the
      gradient, but it changes the filtering distribution and hence the
      posterior—it is a modified model, not a computational trick.
\end{enumerate}

\paragraph{For value-jump targets (integer regime solver).}
\begin{enumerate}
\item \textbf{Reflection/refraction HMC (Section~\ref{sec:event}).} Detect
      regime-path changes, compute the one-sided likelihoods
      $\log p(y\mid\vartheta_b, r^\pm)$, and apply the energy-conserving event
      map \eqref{eq:59}--\eqref{eq:60}. This is exact but expensive: locating
      the implicit boundary requires bisection and full solver + filter
      re-evaluation at intermediate points.
\item \textbf{Discontinuous HMC (Section~\ref{sec:dhmc}).} Use Laplace
      momentum for parameters that affect regime switches. This avoids explicit
      event detection but requires random coordinate order, randomized step
      size, and evaluating the potential on both sides of proposed moves.
\item \textbf{Regime-conditional HMC.} Given a regime path $r_{1:T}$, run
      standard HMC on $p(\vartheta\mid y,r)$ (smooth, fast). Resample $r$ from
      $p(r\mid\vartheta,y)$ using a discrete sampler. Alternate between the two
      steps. Mixing depends on the regime-parameter coupling; if $p(r\mid\vartheta,y)$
      is peaked, regime updates have low acceptance.
\item \textbf{Particle MCMC.} Replace the UKF with a particle filter,
      integrate out the regime path, and use particle Metropolis-Hastings
      (Section~\ref{sec:pmcmc}). This avoids parameter-space discontinuities
      entirely but requires a particle filter run at every proposal.
\end{enumerate}

\subsection{Why DHMC does not solve the continuous-kink problem}\label{sec:dhmc-limitation}

It is tempting to apply discontinuous HMC (Section~\ref{sec:dhmc}) to the UKF
variance-collapse problem, reasoning that ``the target is nonsmooth, and DHMC
is for nonsmooth targets.'' This is incorrect. DHMC with Laplace momentum and
the reflection/refraction event map are designed for \emph{value jumps}: finite
discontinuities in the potential $U(\vartheta)$ itself. The UKF
continuous-kink target has $U$ continuous everywhere; only its gradient jumps.

The DHMC reflection condition checks whether the available kinetic energy
$|p_i|/m_i$ can pay the potential increase $\Delta U$. For a continuous
target, $\Delta U = 0$ at every boundary, so the reflection test is trivially
satisfied, and DHMC degenerates to ordinary coordinate-wise leapfrog with
Laplace momentum. The Laplace momentum may improve mixing in some
high-dimensional settings (constant-velocity motion is cheaper to integrate
than Gaussian trajectories), but it does not address the gradient-jump
inefficiency, because that inefficiency arises from the $O(\epsilon\,\Delta g)$
leapfrog error term, which remains present regardless of the momentum
distribution.

The distinction is sharp: DHMC and reflection/refraction HMC solve the
value-jump problem (Case 3 of Section~\ref{sec:geometries}); they do not solve
the continuous-kink problem (Case 1), which requires either smoother filters,
event-aware gradient switching, or surrogate proposals.

\subsection{Open research questions}\label{sec:ukf-open}

The UKF gradient-discontinuity problem for ZLB models remains only partially
solved. The following questions are active research directions.

\paragraph{Can event-aware integration be made practical?}
Tran and Kleppe's (2025) truncated-step integrator is mathematically valid for
continuous-kink targets, but detecting an implicit OBC boundary requires
bisection with full solver + filter re-evaluation at each bisection step. For
a finite-horizon DSGE model, this can mean solving a 40-period piecewise-linear
system dozens of times per HMC trajectory. Is this cost competitive with
smaller step sizes, or does it dominate the per-effective-sample cost?

\paragraph{Is the marginal likelihood Lipschitz continuous?}
Even with softplus smoothing, the composed map
$\vartheta\mapsto\log p(y\mid\vartheta)$ may fail global Lipschitz continuity.
The filter can amplify small solver changes through covariance updates and log
determinants. Can a local Lipschitz constant be derived as a function of the
softplus temperature, filter type, and time horizon? Or must we accept that
OBC models are inherently non-Lipschitz and design integrators accordingly?

\paragraph{Can neural surrogates learn implicit boundaries?}
A neural score approximation $s_\phi(\vartheta)\approx\nabla_\vartheta\log
p(y\mid\vartheta)$ trained on branch-interior and near-boundary points can
drive HMC proposals with hard-target correction. But implicit OBC boundaries
are high-codimension surfaces defined by solver consistency checks, not simple
hyperplanes. Can a network generalize across such surfaces, or does it
overfit to the training regime and fail when the chain explores a new boundary?

\paragraph{Does UKF variance inflation preserve scientific validity?}
Inflating $P_{t|t-1}$ to prevent collapse smooths the gradient, but it changes
the filtering distribution and hence the posterior. Is there a principled
calibration (e.g., tempering, robust Bayes, conservative updating) under which
the inflated posterior remains a valid answer to the original inference
question? Or is inflation purely a computational heuristic that must be
validated post hoc by comparison with a particle-filter benchmark?

\paragraph{When is regime-conditional HMC competitive?}
Alternating between $p(\vartheta\mid y,r)$ (smooth HMC) and $p(r\mid\vartheta,y)$
(discrete sampler) avoids joint regime-parameter boundaries. But the discrete
sampler can mix poorly when the regime path is long and tightly constrained by
the data. Under what coupling structure does this decomposition outperform
joint methods?

\subsection{Interim recommendation}\label{sec:ukf-recommendation}

Until the research questions of Section~\ref{sec:ukf-open} are resolved, the
safest path for production inference with UKF-filtered ZLB models is:

\begin{enumerate}
\item \textbf{Measure the geometry first.} Run the protocol of
      Section~\ref{sec:ukf-measurement} to classify the target as
      continuous-kink or value-jump and quantify gradient variation.
\item \textbf{For continuous-kink targets:} Use standard HMC with
      aggressively tuned step size, report crossing diagnostics (acceptance
      conditional on steep-gradient crossing), and accept that mixing will be
      slower than for smooth targets. If computational budget allows, test a
      neural surrogate with hard-target correction.
\item \textbf{For value-jump targets:} Standard leapfrog is invalid.
      Reflection/refraction HMC (Section~\ref{sec:event}) is exact but
      expensive. DHMC (Section~\ref{sec:dhmc}) is cheaper but requires random
      coordinate order and step size. Regime-conditional HMC is a fallback if
      event methods prove intractable. Particle MCMC (Section~\ref{sec:pmcmc})
      avoids the discontinuity entirely.
\item \textbf{Validate reversibility.} For any event-aware or DHMC method,
      record a trajectory, flip momentum, replay, and check
      $\|\vartheta_{\text{final}} - \vartheta_{\text{initial}}\| < \epsilon$.
      Reversibility failure indicates a sampler bug or boundary-handling error.
\item \textbf{Report boundary telemetry.} Every UKF-based inference result
      should state: geometry classification (smooth/kink/jump), boundary
      crossing rate, acceptance rate conditional on crossing, and softplus
      temperature (if applicable). Without this telemetry, acceptance rates and
      ESS alone cannot diagnose whether poor mixing is due to the target
      geometry or sampler tuning.
\end{enumerate}

The particle-filter route (Section~\ref{sec:pmcmc}) remains the robust
fallback: it sidesteps the UKF gradient problem by maintaining diversity and
smoothing the parameter-space likelihood. Its limitations—biased gradients in
current differentiable implementations, high computational cost—motivate
continued work on UKF-based methods, but they do not yet justify abandoning
particle methods for production ZLB inference.

\section{Reflection and refraction HMC for piecewise-smooth targets}\label{sec:event}

When the target really does jump, so that the density, not merely its gradient, is
discontinuous across a boundary, leapfrog's energy error at a crossing no
longer vanishes with the step size, and the sampler needs an explicit event
mechanism. The physical analogy is optics: a particle hitting an interface
between media either refracts through it, losing kinetic energy equal to the
potential jump, or reflects off it when it lacks the energy to cross.

\subsection{The event calculation}\label{sec:event-calc}

Let $U$ be smooth inside regions separated by an affine boundary with unit
normal $\nu$. During a free-flight step, decompose the momentum into its
normal and tangential parts,
\begin{equation}\label{eq:57}
 p_\perp=(p^{\mathsf T}\nu)\nu,
 \qquad p_\parallel=p-p_\perp.
\end{equation}
At the boundary, define the potential jump $\Delta U=U(q^+)-U(q^-)$. Energy
conservation across the event requires
\begin{equation}\label{eq:58}
 K(p^+)-K(p^-)=-\Delta U.
\end{equation}
For unit Gaussian mass, if $\|p_\perp\|^2>2\Delta U$ the trajectory has
enough normal kinetic energy to pay the potential toll: it crosses the
boundary with the rescaled normal momentum
\begin{equation}\label{eq:59}
 p_\perp^+=
 \sqrt{\|p_\perp^-\|^2-2\Delta U}\,
 \frac{p_\perp^-}{\|p_\perp^-\|};
\end{equation}
otherwise it reflects,
\begin{equation}\label{eq:60}
 p_\perp^+=-p_\perp^-.
\end{equation}
The tangential momentum is unchanged in both cases. Afshar and Domke (2015,
pp.~2--4) detect the first affine-boundary intersection along the drift,
apply \eqref{eq:59} or \eqref{eq:60}, continue the remaining flight, and
repeat for later boundaries.

One point about volume preservation is easy to get wrong and worth attaching
to the right object. The reflection map \eqref{eq:60} is an orthogonal
reflection with unit absolute Jacobian, but the refraction map \eqref{eq:59}
\emph{in isolation} rescales the normal momentum by $s_-/s_+$ and is not
volume preserving when $\Delta U\neq0$. What Afshar and Domke prove is that
the \emph{assembled} flight--event--flight trajectory step is volume
preserving and reversible: the momentum rescaling at the event is exactly
cancelled by the change in the position leg traversed at the new speed, since
compressing the normal velocity stretches (or shrinks) the set of positions
reached in the remaining flight time by precisely the reciprocal factor.
Combined with leapfrog away from boundaries, this gives detailed balance
after the Metropolis correction (Afshar and Domke 2015, pp.~5--7).

Equations \eqref{eq:57}--\eqref{eq:60} are stated for unit Gaussian mass and
cannot be used verbatim in a preconditioned sampler. Any practical HMC
implementation uses a positive-definite mass matrix $M$ with
$K(p)=\tfrac12 p^{\mathsf T}M^{-1}p$, and the correct decomposition is then
orthogonal in the $M^{-1}$ inner product, not the Euclidean one. For a
boundary with (not necessarily unit) normal $n$, define
\begin{equation}\label{eq:60a}
 a=n^{\mathsf T}M^{-1}n,
 \qquad
 s_-=\frac{n^{\mathsf T}M^{-1}p_-}{\sqrt a}.
\end{equation}
For a potential jump $\Delta U$, the trajectory crosses when
$s_-^2>2\Delta U$, with
\begin{equation}\label{eq:60b}
 s_+=\operatorname{sign}(s_-)\sqrt{s_-^2-2\Delta U},
 \qquad
 p_+=p_-+n\,\frac{s_+-s_-}{\sqrt a},
\end{equation}
and reflects otherwise with
\begin{equation}\label{eq:60c}
 p_+=p_--2n\,\frac{n^{\mathsf T}M^{-1}p_-}{n^{\mathsf T}M^{-1}n}.
\end{equation}
Both maps change only the momentum component along $M^{-1}n$, conserve
$K(p)+U(q)$ across the event, and reduce to \eqref{eq:59}--\eqref{eq:60}
when $M=I$ and $\|n\|=1$. Their Jacobians differ: the reflection
\eqref{eq:60c} is an $M^{-1}$-orthogonal reflection with unit absolute
Jacobian, while the refraction \eqref{eq:60b} alone has absolute Jacobian
$|s_-/s_+|\neq1$ for $\Delta U\neq0$, and volume preservation again holds
only for the composed flight--event--flight step, for the same
cancellation reason as above. A correctness test must therefore target the
composed trajectory map (or account for the event Jacobian explicitly in the
acceptance ratio), never the isolated refraction. An implementation must
additionally fix the boundary orientation convention (the sign of $n$ and
hence of $\Delta U$) and the treatment of exact-equality and grazing events
($s_-^2=2\Delta U$ or $s_-=0$), and should be exercised with dense and
diagonal nonidentity metrics, not only with $M=I$.

\subsection{What transfers to a ZLB}\label{sec:event-zlb}

This method is a good match when the posterior in the sampled variables is
piecewise smooth and the switching surfaces are explicit affine boundaries,
or at least surfaces whose first intersection with a trajectory can be found
reliably. A nonlinear endogenous regime solver raises the bar. In that case
the boundary is an implicit surface
\begin{equation}\label{eq:61}
 b(q)=0,
\end{equation}
and a trajectory needs (i) a first-root solver for $b(q(t))$, (ii) the
correct one-sided values of the full log posterior, and (iii) an argument
that the selected equilibrium branch is single-valued at the event. The
Metropolis correction absorbs integration error relative to a well-defined
target; when a root search fails or the solver lands on a different fixed
point, the two one-sided posterior values in \eqref{eq:5b} are no longer
defined, so the problem statement itself, and not the sampler, is what needs
repair.

\section{Discontinuous HMC with Laplace momentum}\label{sec:dhmc}

Reflection--refraction HMC pays for every boundary crossing individually.
When a discrete or ordinal variable is embedded into the continuous state,
the trajectory may cross hundreds of density jumps per proposal, and a
different momentum distribution turns out to make each crossing dramatically
cheaper.

\subsection{Embedding an ordinal regime}\label{sec:dhmc-embed}

Nishimura, Dunson, and Lu (2020, pp.~366--368) embed an integer parameter $N$
in a real variable $\widetilde N$:
\begin{equation}\label{eq:62}
 N=n \iff \widetilde N\in(a_n,a_{n+1}],\qquad
 \widetilde\pi(\widetilde n)=
 \sum_n\frac{\pi_N(n)}{a_{n+1}-a_n}
 \mathbf 1\{a_n<\widetilde n\leq a_{n+1}\}.
\end{equation}
The interval length in the denominator is essential: it is the Jacobian-like
factor that makes the total embedded mass of each interval equal to
$\pi_N(n)$, so that integrating $\widetilde\pi$ over $(a_n,a_{n+1}]$ recovers
the original probability. The embedding also imposes an ordering on the
discrete values; an arbitrary ordering can place high-probability values in
separated intervals, and energy-conserving dynamics may then mix poorly
between the resulting modes.

\subsection{Why Gaussian momentum is expensive, and the Laplace alternative}\label{sec:dhmc-laplace}

When a Gaussian-momentum trajectory crosses many density jumps, every
crossing must be located and corrected as in Section~\ref{sec:event-calc}.
For a discrete coordinate with posterior uncertainty of hundreds of units,
that means hundreds of target evaluations per trajectory. Nor can one simply
ignore the jumps and let the Metropolis step clean up: the leapfrog error at
a jump does not vanish as $\epsilon\to0$, because the unaccounted potential
change is finite regardless of the step size (Nishimura, Dunson, and Lu 2020,
pp.~367--370).

Choose instead independent Laplace momentum,
\begin{equation}\label{eq:63}
 K(p)=\sum_i |p_i|/m_i,
\qquad
 \dot q_i=m_i^{-1}\operatorname{sign}(p_i).
\end{equation}
The velocity now depends only on the sign of the momentum, not its
magnitude, so the position update over a step is known in advance regardless
of how much kinetic energy the coordinate loses along the way. For
coordinate $i$, propose the endpoint
\begin{equation}\label{eq:64}
 q_i^\star=q_i+\epsilon m_i^{-1}\operatorname{sign}(p_i),
\qquad
 \Delta U=U(q^\star)-U(q).
\end{equation}
If $|p_i|/m_i>\Delta U$, accept the endpoint and reduce the momentum by the
potential increase,
\begin{equation}\label{eq:65}
 p_i\leftarrow p_i-\operatorname{sign}(p_i)m_i\Delta U;
\end{equation}
otherwise keep $q_i$ and reflect the momentum, $p_i\leftarrow-p_i$. The
coordinate maps are applied sequentially in a randomly permuted order; the
random permutation is needed for reversibility in distribution. Nishimura
et al.\ prove that the coordinate map is volume preserving and reversible
almost everywhere, and that the assembled integrator is valid with a final
Metropolis correction (Nishimura, Dunson, and Lu 2020, pp.~371--374). They
also randomize the step size to avoid a fixed-grid reducibility problem.

The algorithm is attractive for an explicit regime coordinate because a
single endpoint likelihood evaluation in \eqref{eq:64} can cross several
ordinal intervals at once; this is exactly the economy that Gaussian
momentum lacks. It is a poor fit for a solver that hides a complicated
regime path inside one black-box function: the endpoint energy difference
$\Delta U$ is meaningful only if both endpoints evaluate the same
well-defined posterior, the density of the model actually being
estimated, and all branch ambiguity in the solver is resolved.

\section{Mixed HMC for an explicit regime path}\label{sec:mixed}

Embedding a discrete variable into the real line, as in
Section~\ref{sec:dhmc-embed}, forces an ordering that may not exist
naturally. Zhou's mixed HMC (2020, pp.~2--4) instead retains the discrete
variable as genuinely discrete. Let $z$ be discrete, $q$ continuous, and
$\pi(z,q)\propto e^{-U(z,q)}$. Attach to each discrete coordinate $j$ a
clock $s_j$ living on a flat torus and a discrete momentum $p_j^D$. The
continuous coordinates evolve by ordinary Hamiltonian dynamics; when a clock
reaches its boundary, the sampler proposes $\widetilde z$ from an
irreducible single-site proposal $Q_j(\widetilde z\mid z)$ and treats the
move as an event. Define
\begin{equation}\label{eq:66}
 \Delta E=\log\frac{\pi(z,q)Q_j(\widetilde z\mid z)}
                         {\pi(\widetilde z,q)Q_j(z\mid\widetilde z)}.
\end{equation}
For a kinetic energy $k_j^D(p_j^D)$, the discrete move is accepted (the
trajectory ``crosses'') with the momentum changed so that
\begin{equation}\label{eq:67}
 k_j^D(p_j^{D,+})=k_j^D(p_j^{D,-})-\Delta E
\end{equation}
when the right-hand side is nonnegative; otherwise the momentum reflects
and $z$ is kept. This is the refraction--reflection rule of
Section~\ref{sec:event-calc} transplanted to a discrete jump, with the
proposal-corrected energy difference \eqref{eq:66} playing the role of
$\Delta U$. The continuous coordinates are advanced by a reversible
volume-preserving integrator, and a final Metropolis correction gives
invariance of the joint distribution. Note that the theorem depends on the
proposal ratio appearing in \eqref{eq:66}, not on a conditional Gibbs draw.

For a \emph{stochastic} ZLB path, one can set $z=r_{1:T}$, the full regime
path, or use a blocked path proposal. Equation \eqref{eq:66} then includes
the complete regime-path probability and the nonlinear state-space
likelihood. This is conceptually clean but can be expensive: changing one
regime indicator may require re-solving a long state path to evaluate the
new likelihood. For the deterministic indicator of \eqref{eq:5a}, by
contrast, the regime path is a function of the continuous state, a
regime-only proposal has zero support, and this construction does not apply
at all. Either way, the method presupposes an explicit probability law over
regime paths, including a stance on multiple equilibrium paths, before
the energy difference in \eqref{eq:66} is even defined.

\section{Smooth categorical relaxations and nonsmooth energies}\label{sec:relaxations}

If exact discrete dynamics are too expensive, one can instead smooth the
discreteness away and run ordinary HMC on the relaxed target, accepting
that the relaxed posterior is not the original one. The Concrete /
Gumbel-softmax method (Torgander, Magnusson, and Wallin 2024, pp.~2--3) maps
categorical probabilities to a continuous simplex variable. As the
temperature $\tau\to0$, the variable becomes nearly one-hot, but the
transformed posterior develops steep geometry: gradients with respect to
simplex coordinates can diverge near zero, and the temperature direction
becomes ill-conditioned. It is therefore a useful relaxation baseline, not
an exact categorical or ZLB sampler, and any comparison against it should
report the temperature used and the induced posterior discrepancy.

A different smoothing tradition targets nonsmooth but continuous energies.
Chaari et al.'s nonsmooth HMC (2016, pp.~1--5) uses subgradients or a
proximity operator
\begin{equation}\label{eq:68}
 \operatorname{prox}_{\phi}(x)=
 \arg\min_y\left\{\phi(y)+\tfrac12\|y-x\|^2\right\}
\end{equation}
for convex nondifferentiable energies such as an $\ell_1$ penalty. The
distinction of Section~\ref{sec:kink-validity} applies here with full force:
the proximity operator \eqref{eq:68} is built for convex \emph{continuous}
energies, whose nondifferentiability is a kink, whereas a density
discontinuity is a finite jump in $U$ that no subgradient or proximal
smoothing removes, and an equilibrium regime selection is a modeling choice,
not an energy. Prox-HMC therefore enters the ZLB problem exactly when one
deliberately chooses a convex nonsmooth approximation of the constrained
model.

\section{Particle MCMC, pseudo-marginal HMC, and differentiable state-space HMC}\label{sec:pmcmc}

Every gradient method above struggles with the same obstacle: at
the zero lower bound the exact likelihood is an integral over latent paths
whose integrand is only piecewise smooth. Particle methods attack the problem
from the opposite direction. Rather than smoothing the model, they replace the
intractable likelihood by an unbiased Monte Carlo estimate and then build a
Markov chain that is \emph{exactly} invariant for the posterior despite the
estimation noise. What follows develops that construction, asks when it can
be combined with Hamiltonian dynamics, and examines what modern
``differentiable'' particle filters do and do not deliver.

\subsection{Particle MCMC's exact extended target}\label{sec:pm-target}

One disarmingly simple identity carries the whole construction. Let $u$ collect \emph{all} random variables
consumed by a run of the particle filter, every proposal draw and every
resampling variable, with joint density $m_\vartheta(u)$, and let
$\widehat L(\vartheta,u)\geq 0$ be the resulting likelihood estimate. The
standard particle filter is unbiased:
\begin{equation}\label{eq:69}
 \int \widehat L(\vartheta,u)\,m_\vartheta(u)\,du = L(\vartheta).
\end{equation}
Now consider the \emph{extended} target on the pair $(\vartheta,u)$,
\begin{equation}\label{eq:70}
 \bar\pi(\vartheta,u) \propto p(\vartheta)\,\widehat L(\vartheta,u)\,m_\vartheta(u).
\end{equation}
Integrating \eqref{eq:70} over $u$ and applying \eqref{eq:69} shows that its
$\vartheta$-marginal is exactly $p(\vartheta)L(\vartheta)$: the posterior we
actually want, with no bias from the finite particle count. This is the
pseudo-marginal identity underlying particle MCMC (Andrieu, Doucet, and
Holenstein 2010). Its practical force is that a Metropolis--Hastings chain
targeting \eqref{eq:70}, in which each proposed $\vartheta'$ comes with a fresh
filter run $u'$ and the acceptance ratio uses the \emph{estimated} likelihoods,
is a valid MCMC algorithm for the exact posterior. The estimation noise is
part of an enlarged state space on which the chain is exact, not an
approximation error left over in the parameter marginal.

\subsection{What pseudo-marginal HMC additionally requires}\label{sec:pmhmc}

Given \eqref{eq:70}, it is tempting to go further: if we can evaluate
$\log\widehat L(\vartheta,u)$, why not differentiate it and run HMC on the
extended target? Alenl\"ov, Doucet, and Lindsten (2021, Secs.~2--3) construct
exactly such an extended Hamiltonian, and their analysis makes precise what the
shortcut demands. The construction requires, explicitly:
\begin{enumerate}
\item a nonnegative unbiased estimator $\widehat L(\vartheta,U)$;
\item auxiliary variables $U\sim N(0,I)$;
\item continuous differentiability of the map
      $(\vartheta,U)\mapsto\widehat L(\vartheta,U)$; and
\item pointwise evaluability of $\nabla\log\widehat L(\vartheta,U)$.
\end{enumerate}
Under these conditions the algorithm is exact for the extended target, and as
the number of importance samples grows the trajectories approach those of ideal
HMC on the marginal. The catch is condition (iii). Alenl\"ov, Doucet, and
Lindsten note explicitly that ordinary particle-filter likelihoods typically
\emph{fail} it: a multinomial resampling index changes discontinuously when a
uniform random number crosses a cumulative-weight threshold, so
$\widehat L$ is discontinuous in $(\vartheta,U)$ even when the underlying model
is smooth.

The consequence for ZLB inference is sharp. An unbiased bootstrap particle
filter supports particle Metropolis--Hastings via the identity of
Section~\ref{sec:pm-target}, while pseudo-marginal HMC additionally needs
condition (iii), which resampling breaks. A differentiable resampling scheme
or a transport-based relaxation can restore
condition (iii), but the repair is not free: the gradient of a smoothed or
transport-resampled computation is either an approximation, or the exact
gradient of a \emph{different} extended target; and it remains the latter
only once that target's full law and Metropolis correction have been written
down. Pseudo-marginal HMC therefore rests on a separately constructed
differentiable extended estimator, a genuine design task beyond a working
particle Metropolis--Hastings sampler. The correlated pseudo-marginal method
(Deligiannidis, Doucet, and Pitt
2018) addresses a different axis: it improves particle Metropolis--Hastings
mixing by correlating
successive $U$ draws. That makes it the right variance-reduction comparator,
while the differentiability repair that pseudo-marginal HMC needs must come
from the resampling scheme itself.

\subsection{Joint latent-state HMC}\label{sec:joint-hmc}

Instead of marginalizing the states by particle filtering, one can sample
$(\vartheta,x_0,\epsilon_{1:T})$ directly, using the joint density
\eqref{eq:56}. This avoids a particle likelihood altogether and works well for
smooth nonlinear models (Childers et al.\ 2022, pp.~10--16). At the ZLB,
however, the structure of the joint density changes with the regime treatment.
With a stochastic explicit regime path, \eqref{eq:56} becomes a mixed target:
the state transition and observation density are piecewise in $r_t$, and the
discrete path needs a transition kernel that actually has support on regime
changes. With a deterministic regime, the density instead concentrates on the
support described in \eqref{eq:5a}, and the discontinuity migrates into the
continuous coordinates. Joint HMC thus supplies the continuous-coordinate
component of a complete sampler for this problem class; the regime treatment
determines what has to be composed with it.

\subsection{Particle Gibbs and ancestor sampling}\label{sec:pgas}

Particle Gibbs takes yet another route: instead of estimating the marginal
likelihood, it samples an entire latent trajectory as a Gibbs block. Given a
reference path $x'_{1:T}$, conditional sequential Monte Carlo runs a particle
filter in which one particle is pinned to that path. Plain particle Gibbs has a
well-known pathology: when the particle genealogy collapses, the sampled path
tends to retain its early ancestors, so mixing over the early part of the
trajectory stalls. Particle Gibbs with ancestor sampling (PGAS) repairs this by
resampling the fixed particle's ancestor at each time step (Lindsten, Jordan,
and Sch\"on 2014, Secs.~3--5).

For a Markov state-space model, one PGAS sweep is implementable as follows. At
$t=1$, draw particles $1{:}M-1$ from the initial proposal, set $X_1^M=x'_1$,
and weight all $M$ particles. For each $t=2{:}T$:
\begin{enumerate}
\item For $j<M$, draw $A_t^j\sim\operatorname{Categorical}(W_{t-1})$ and
      propagate $X_t^j$ from the proposal in \eqref{eq:12b}.
\item Set $X_t^M=x'_t$, but draw its ancestor with probabilities
\begin{equation}\label{eq:70a}
 \Pr(A_t^M=j)\propto W_{t-1}^j\,
 p(x'_t\mid X_{t-1}^j,\vartheta).
\end{equation}
\item Recompute the importance weights. At $T$, draw a terminal particle from
      $W_T$ and trace its ancestors to obtain the new path.
\end{enumerate}
Equation \eqref{eq:70a} is the Markov simplification of the general ancestor
weight, which multiplies the old particle weight by the ratio of the target
density after and before attaching the fixed future to a candidate past.
Lindsten, Jordan, and Sch\"on prove that this kernel leaves the joint smoothing
distribution invariant for \emph{every} particle count; their proof interprets
the sweep as a partially collapsed Gibbs update on an extended particle system
(Lindsten, Jordan, and Sch\"on 2014, Algorithm~2 and Theorem~1).

How does this apply at the bound? For a deterministic ZLB, the particle state
is the continuous state or shock path, and the regime $r_t=R(q)$ is derived
from it; PGAS then mixes over paths, and the regime sequence follows. For a
stochastic regime, one takes $(r_t,x_t)$ as the particle state, and PGAS mixes
over regimes and states jointly. What PGAS does \emph{not} license is an
independent update of a deterministic, zero-support regime label: if the regime
is a function of the continuous path, there is no separate coordinate to
resample.

\subsection{What differentiable particle filters do and do not supply}\label{sec:dpf}

Two recent lines of work make particle filters interact with automatic
differentiation, and each delivers something precise but different.

\'Scibior and Wood keep the ordinary particle filter's forward values entirely
unchanged but insert stop-gradient weight ratios, so that automatic
differentiation of the resulting computation returns the Fisher-identity score
estimator associated with the particle genealogy (\'Scibior and Wood 2021,
Algorithm~1 and Sec.~3). Poyiadjis, Doucet, and Singh derive the underlying
path-space and forward-smoothing score estimators and show their different cost
and variance growth (Poyiadjis, Doucet, and Singh 2011, Secs.~2--3). The
estimator is consistent as the particle count grows. At any finite particle
count, however, the score carries Monte Carlo noise, so feeding it to a
leapfrog integrator as if it were the deterministic force of the marginal
log likelihood produces a stochastic-gradient chain whose invariant
distribution is no longer covered by the usual HMC exactness proof.

Corenflos et al.\ instead replace discrete resampling with a differentiable
entropy-regularized optimal-transport map. This changes the finite-particle
forward algorithm itself and introduces a regularization parameter; their
consistency analysis requires the particle count to grow \emph{and} the
regularization to shrink (Corenflos et al.\ 2021, Sec.~4).

Both constructions are important comparison methods. Exact use of either
inside HMC requires the full pseudo-marginal discipline of
Section~\ref{sec:pmhmc}: an explicitly written extended target, random numbers held
fixed during each reversible trajectory, and a Metropolis energy computed for
that same extended density. Inserting either gradient as a noisy force into
ordinary HMC skips all three steps and with them the exactness proof.

\section{What the literature supports}\label{sec:synthesis}

Standing back from the individual methods, the literature surveyed so far
supports a layered architecture for anyone designing Bayesian inference at an
occasionally binding bound. The four layers are connected but have different
jobs, and much practical grief comes from asking one layer to do another's
work.

\textbf{Layer A: a verified linear benchmark.} Implement OccBin-style piecewise
path iteration and the conditional Kalman likelihood \eqref{eq:8}--\eqref{eq:11},
and validate them against a published example and against synthetic data with
known regime paths. This layer looks pedestrian, but it is where timing
conventions, inequality signs, likelihood normalization, missing observations,
and multiple-fixed-point diagnostics get established; errors here propagate
invisibly into every later layer.

\textbf{Layer B: a nonlinear likelihood authority.} Implement the bootstrap
particle filter \eqref{eq:12a}--\eqref{eq:13} and, where the decision rules are
piecewise linear and continuous, the Aruoba et al.\ conditionally optimal
special case. The bootstrap filter handles any simulable nonlinear transition;
the conditionally optimal filter is the low-variance bridge for
piecewise-linear continuous decision rules. Particle Metropolis--Hastings then
supplies, via \eqref{eq:70}, a finite-particle extended target whose parameter
marginal is the exact posterior: the ``ground truth'' sampler against which
faster methods are judged.

\textbf{Layer C: deterministic complementarity.} Here the regime is the derived
quantity $R(q)$, not an independent unknown. The safe starting point is
joint-state random-walk or slice updates and particle MCMC. Event-aware HMC
should be added only after the complete nonlinear solution map supplies a
first-boundary oracle and both one-sided posterior values at a crossing.
Within one region, a fixed-branch HMC update explores efficiently; crossing
into the other regime, however, requires moving the continuous state across
the boundary, because by \eqref{eq:5a} a proposal that changes only the
regime label has zero target density.

\textbf{Layer D: stochastic or explicitly selected regimes.} When both regime
values carry positive probability, alternate fixed-path HMC with PGAS or
another valid discrete Metropolis kernel. Zhou's mixed HMC becomes an
acceleration candidate only after this simpler composition has been checked
against exact enumeration on a toy model: its clock changes update timing, not
the target distribution, so agreement on a small exactly solvable case is the
right admission test.

Throughout, the Concrete-relaxation and softplus routes of
Section~\ref{sec:relaxations} should be understood as approximation arms. In
particular, the shadow-rate case study of Section~\ref{sec:shadow} uses a smooth
softplus bound and is therefore a useful comparator for the hard-bound methods,
not evidence about genuinely discontinuous semantics.

The two case studies later in this survey instantiate these layers in
complementary ways. The two-country shadow-rate model of Section~\ref{sec:shadow} is a
hard-bound observation-map kink target: Layer B specializes to the bootstrap
particle-filter likelihood of Section~\ref{sec:shadow-filter}, with the ideal cell
decomposition \eqref{eq:84}--\eqref{eq:87} as its proposal and diagnostic
layer; Layer C specializes to joint kinked-target HMC under the validity
conditions of Section~\ref{sec:kink-validity}; and Layers A and D do not apply to that
model as specified. The New Keynesian model of Section~\ref{sec:nk} is the converse:
Layers A, C, and D carry the weight, the genuine discontinuity lives in the
solution correspondence and its selection rule (Section~\ref{sec:nk-mult}), and Layer D
becomes available exactly when selection is made stochastic
(Section~\ref{sec:nk-sunspot}).

\section{Exact targets and implementable transition kernels}\label{sec:targets}

A method is only as well specified as the object it targets, so we
now write those objects down. For each of the three regime treatments (deterministic
threshold, stochastic regime, and multiple equilibria) we state the exact
posterior density and a transition kernel that provably preserves it. Being
explicit here pays off later: many subtle errors in bound-aware samplers are
errors about \emph{which} distribution is being sampled, not about the
dynamics.

\subsection{Deterministic threshold target}\label{sec:target-det}

Let $q=(\vartheta,x_0,\epsilon_{1:T})$, recursively define
$r_t=\mathcal R(x_{t-1},\epsilon_t,\vartheta)$ and
$x_t=F_{r_t}(x_{t-1},\epsilon_t,\vartheta)$, and assume this solution is
unique. Because the regime is a deterministic function of $q$, the posterior is
a density on the continuous variables alone:
\begin{equation}\label{eq:71}
 \Pi_D(q)\propto p(\vartheta)\,p(x_0\mid\vartheta)
 \prod_{t=1}^T p(\epsilon_t\mid\vartheta)\,
 g_{r_t(q)}\bigl(y_t\mid x_t(q),\vartheta\bigr).
\end{equation}

An event-aware HMC step for \eqref{eq:71} proceeds as follows. Inside a region
with fixed $r_{1:T}$, differentiate the same finite program that evaluates
\eqref{eq:71}: the gradient of the density actually being computed, not of a
surrogate. At the first event $b_k(q(t))=0$ along the trajectory, evaluate the
two one-sided values of the \emph{complete} posterior, compute the potential
jump \eqref{eq:5b}, and apply the general-mass event maps \eqref{eq:60b} or
\eqref{eq:60c}; the simpler maps \eqref{eq:59}/\eqref{eq:60} apply only in the
unit-mass special case. Continue the flight for the unused part of the position
step. At the end of the trajectory, reverse the momentum and accept with
\begin{equation}\label{eq:72}
 \alpha_D=1\wedge\exp\{H(q,p)-H(q',p')\}.
\end{equation}

This full step is implementable only under real conditions: the event locator
must return the \emph{earliest} crossing; the branch solver must be
deterministic and single-valued on each side of the boundary; and the
\emph{composed} flight--event--flight trajectory map must be reversible and
volume preserving (or the event Jacobian must be accounted for explicitly in
the acceptance ratio). Since the isolated refraction map rescales the normal
momentum (Section~\ref{sec:event-calc}), volume preservation holds only for
the composed map, so the composition is what a correctness test must
exercise. If any of these conditions is missing, particle
Metropolis--Hastings using \eqref{eq:13}, or a non-gradient joint-state kernel,
remains the exact baseline.

\subsection{Stochastic mixed target}\label{sec:target-stoch}

If regimes carry genuine transition mass, the regime path is an unknown in its
own right and the joint target is instead
\begin{equation}\label{eq:73}
 \Pi_S(\vartheta,x_{0:T},r_{1:T})\propto
 p(\vartheta)\,p(x_0\mid\vartheta)
 \prod_{t=1}^T p(r_t\mid r_{t-1},x_{t-1},\vartheta)\,
 p(x_t\mid x_{t-1},r_t,\vartheta)\,
 g_{r_t}(y_t\mid x_t,\vartheta).
\end{equation}

A transparent first kernel composes two invariant updates:
\begin{enumerate}
\item Hold $r_{1:T}$ fixed and apply ordinary Metropolized HMC to
      $(\vartheta,x_{0:T})$, using the gradient of $\log\Pi_S$. With the
      regime path frozen, the target is smooth in the continuous coordinates,
      so standard HMC theory applies.
\item Hold the continuous variables fixed and update one regime, a regime
      block, or the whole latent path with Metropolis, forward-filter
      backward-sampling, or PGAS. For a proposal $Q(r'\mid r,q)$, a
      regime-only move is accepted with
\begin{equation}\label{eq:74}
 \alpha_R=1\wedge
 \frac{\Pi_S(q,r')\,Q(r\mid r',q)}
      {\Pi_S(q,r)\,Q(r'\mid r,q)}.
\end{equation}
\end{enumerate}
Each component preserves $\Pi_S$, so their composition does as well; this is
the standard Metropolis-within-Gibbs argument, and it requires no smoothness
across regime changes. Mixed HMC replaces the second component by the clock
dynamics \eqref{eq:66}--\eqref{eq:67} and interleaves it with the first. One
refinement deserves flagging: if a regime proposal jointly transforms the
continuous coordinates, say to keep an economic constraint satisfied after
the regime flips, then the move is no longer discrete-only, and its reverse
proposal density and the absolute Jacobian of the continuous transformation
must be added to \eqref{eq:74}.

\subsection{Multiple equilibria}\label{sec:target-multi}

Finally, suppose the model admits multiple solutions
$s_t\in\mathcal S(x_{t-1},\epsilon_t,\vartheta)$ at some states, as
occasionally binding constraint models genuinely can. Then one must specify a
selection mass $p(s_t\mid x_{t-1},\epsilon_t,\vartheta)$ and either include
$s_t$ as a latent variable in \eqref{eq:73} or integrate it out. Without that
mass, neither \eqref{eq:71} nor \eqref{eq:73} is a defined target: the
likelihood of the data is simply not determined by the structural parameters.
Supplying the selection mass is therefore part of the economic model
specification; the sampler's job begins only once the target it must draw
from exists.

The two case studies that follow instantiate these abstract targets:
Section~\ref{sec:shadow} develops a two-country shadow-rate term-structure model whose
bound enters through the observation map, and Section~\ref{sec:nk} develops a small New
Keynesian model whose bound enters through the solution correspondence itself.

\section{Case study I: a two-country shadow-rate term-structure model}
\label{sec:shadow}

Near a lower bound on nominal rates, observed yield curves flatten and
compress against the bound while the economically meaningful latent
state, the ``shadow'' curve that would prevail in the absence of the
bound, keeps moving. A yield pinned at the bound for two years is
consistent with a shadow rate slightly below it or far below it, and
distinguishing those two situations is exactly what monetary-policy
analysis at the bound requires. We therefore want Bayesian inference on
the latent shadow factors and the model parameters from bounded yield
observations, and this is a state-space problem in which the bound
enters the observation map as a censoring nonlinearity. This case study
builds a concrete two-country shadow-rate term-structure model from
scratch, places it in the literature, and then applies the machinery of
the preceding sections: which posterior geometry the model has, which
filters and samplers are exact for which version of it, and where the
smooth approximation that makes gradient-based sampling easy stops
being innocent.

\subsection{The model}
\label{sec:shadow-model}

The latent state is
\(x_t=(x_t^{d\mathsf T},x_t^{f\mathsf T},x_t^{b\mathsf T})^{\mathsf T}
\in\mathbb R^8\): domestic (\(d\)) and foreign (\(f\)) dynamic
Nelson--Siegel (DNS) factor triples (level, slope, curvature) and
two directed FX-basis factors. The transition is a stationary Gaussian
VAR(1) toward a long-run mean,
\begin{equation}\label{eq:75}
 x_t=(I-\Phi)\bar\theta+\Phi x_{t-1}+\eta_t,
 \qquad \eta_t\sim N(0,Q),
\end{equation}
and in the working example below the estimated parameters \(\vartheta\)
are components of \(\bar\theta\) together with grouped log
measurement-noise scales, while \(\Phi\), \(Q\), and the initial
covariance are held fixed: a deliberately small nine-parameter chart
that keeps attention on the bound rather than on the VAR dynamics. The
shadow instantaneous forward curve of country \(c\in\{d,f\}\) at horizon
\(s\) is affine in that country's factors with a fixed decay
\(\lambda_c\):
\begin{equation}\label{eq:76}
 f_c(s;x)=a_c(s)^{\mathsf T}x^c,
 \qquad
 a_c(s)=\bigl(1,\;e^{-\lambda_c s},\;\lambda_c s\,e^{-\lambda_c s}
 \bigr)^{\mathsf T}.
\end{equation}

The bound is imposed on the \emph{forward} curve, not on yields
directly. This is the Krippner-style architectural choice, and
Section~\ref{sec:shadow-lit} explains why it matters; for now, the point is
that a yield is a maturity average of forwards, so bounding forwards
and then averaging produces a different, and economically more
natural, object than bounding the average. The bound enters through
one of two scalar maps applied to the forward curve, the hard censoring
map and its softplus smoothing,
\begin{equation}\label{eq:77}
 m_\ell(u)=\max\{\ell,u\},
 \qquad
 s_{\ell,\alpha}(u)=\ell+\alpha\log\!\left(1+e^{(u-\ell)/\alpha}\right),
 \qquad \alpha>0.
\end{equation}
The working example uses \(s_{\ell,\alpha}\) exclusively, with
per-country constants
\((\lambda_d,\ell_d,\alpha_d)=(0.65,\,0,\,1.5\times10^{-3})\) and
\((\lambda_f,\ell_f,\alpha_f)=(0.45,\,-0.005,\,1.0\times10^{-3})\); the
foreign bound is a negative effective lower bound, and \(\alpha\) is a
fixed model constant, not an estimated parameter. Observed zero-coupon
yields are maturity averages of the bounded forward curve. The integral
has no closed form once the softplus (or the max) is inside it, so it
is evaluated by an order-\(K\) Gauss--Legendre rule on \([0,1]\) (the
working example fixes \(K=40\)):
\begin{equation}\label{eq:78}
 y_c(\tau_i;x)=\frac1{\tau_i}\int_0^{\tau_i}
 s_{\ell_c,\alpha_c}\bigl(f_c(u;x)\bigr)\,du
 \;\approx\;\sum_{k=1}^{K}w_k\,
 s_{\ell_c,\alpha_c}\bigl(f_c(\tau_i v_k;x)\bigr),
\end{equation}
with six maturities per country. Quadrature is used because the
integrand is smooth and cheap, and a fixed high-order rule keeps the
observation map deterministic in \(x\), an essential property once
we start differentiating through it. The sole FX observation is the raw
log forward with observed log spot subtracted, tied to the two yield
curves and the basis factors by the covered-interest-parity identity
\begin{equation}\label{eq:79}
 \log F_t(\tau)-\log S_t=\tau\bigl(
 y_d(\tau;x_t)-y_f(\tau;x_t)+b(\tau;x_t^b)\bigr),
\end{equation}
where \(b(\tau;\cdot)\) is affine in the basis factors. All measurement
noise is diagonal Gaussian with grouped scales.

In the implementation we studied, the active inference route evaluates
the marginal likelihood of \eqref{eq:75}--\eqref{eq:79} with a
square-root unscented Kalman filter and runs HMC and NeuTra-HMC over
the nine-parameter chart against that filter's value and score. Two
facts about that route follow immediately, and they frame everything
in this case study. First, the resulting posterior is a
\emph{smooth-model} posterior: \(s_{\ell,\alpha}\) is \(C^\infty\), so
no discontinuity of any kind is present in the model actually being
estimated. Second, the UKF likelihood is itself a deterministic,
well-defined function of the parameters, so HMC on it is exact
\emph{for that target}; the open question is its distance from the
exact-model likelihood, which Section~\ref{sec:shadow-filter} and
Section~\ref{sec:shadow-routes} address.

\subsection{Where this sits in the shadow-rate literature}
\label{sec:shadow-lit}

Where does this design sit in the shadow-rate term-structure family? It is a
reduced-form member of that family, and the survey's earlier machinery maps
onto the family precisely. Black (1995) introduced the observation that currency makes
the nominal short rate an option, \(r_t=\max\{0,r_t^\star\}\) (we cite
this for orientation; we were unable to obtain the full text). Krippner
(2012, Secs.~1--2) made the multi-factor Gaussian case tractable by
bounding the \emph{forward} curve (the ZLB forward is the shadow
forward plus a closed-form currency-option effect) and obtaining
yields by elementary numerical integration of the bounded forward
curve; our working example follows exactly this
bounded-forward-then-integrate architecture with the softplus in place
of Krippner's Gaussian-based map. Priebsch (2013) developed the
cumulant-based arbitrage-free approximation for the same class. Wu and
Xia (2016, Sec.~2.3, eqs.~(6)--(8)) write the one-period forward
observation as
\(\underline r+\sigma^Q\,g\!\left((a_n+b_n^{\mathsf T}X_t-\underline r)/
\sigma^Q\right)\) with the closed-form \(g(z)=z\Phi(z)+\phi(z)\) and
estimate by extended Kalman filtering, noting that monotonicity of
\(g\) keeps the likelihood surface well behaved. Christensen and
Rudebusch (2015, Appendix~B) estimate arbitrage-free shadow-rate
Nelson--Siegel models with the extended Kalman filter built on the
discretized Ornstein--Uhlenbeck transition, literally the
\((I-e^{-K\Delta})\theta+e^{-K\Delta}X_{t-1}\) form of \eqref{eq:75}.
Kim and Priebsch (2013, Sec.~4, eqs.~(8)--(9)) replace the extended
filter with the \emph{unscented} Kalman filter for
quasi-maximum-likelihood estimation, documenting that first-order
linearization can be numerically unstable; every earlier zero-bound
study they name (Ichiue--Ueno, Kim--Singleton, and the 2013
Christensen--Rudebusch draft) had used the extended filter. Lemke
and Vladu (2017) estimate a euro-area shadow-rate model with a
time-varying, estimated lower bound using the extended filter. Finally,
Opschoor and van der Wel (2024, Secs.~2.1--2.3, eqs.~(2)--(7)) build
the reduced-form smooth shadow-rate DNS model closest to our working
example: AR(1) DNS factors, a lower bound applied through a smooth
approximation of the max (their menu is the Gaussian-based map
\(f_G(z)=z\Phi(z)+\phi(z)\), the softplus \(f_S\), and an
inverse-exponential map, each with a sharpness scaling) and two-step
least-squares estimation, with nonlinear state-space estimation via the
extended Kalman filter noted as the alternative.

Three points of contrast matter for what follows. First, Opschoor and
van der Wel bound the \emph{yield} directly,
\(y^Z(\tau)=r_{LB}+f(y(\tau)-r_{LB})\), whereas Krippner-type models,
our working example among them, bound the \emph{forward} curve and
integrate. The two choices differ: by pointwise convexity,
\begin{equation}\label{eq:80}
 \sum_{k}w_k\max\{\ell,\,a(s_k)^{\mathsf T}x\}
 \;\geq\;
 \max\Bigl\{\ell,\;\sum_{k}w_k\,a(s_k)^{\mathsf T}x\Bigr\},
\end{equation}
with equality only when the forward curve does not cross the bound
inside the maturity, so forward-level censoring produces weakly higher
yields near the bound and the two models are not reparameterizations of
each other. Second, Opschoor and van der Wel defend smoothing as an
\emph{economic} specification (yields empirically leave the bound
gradually), which licenses reading \(\alpha\) as a structural
smooth-transition constant rather than as a numerical tolerance;
Section~\ref{sec:shadow-softplus} keeps both readings apart. Third, the
domain-standard estimators in this line are all single-Gaussian
moment-closure filters (extended or unscented), exactly the object
whose branch-mass deficiency Section~\ref{sec:ukf-branchloss} quantifies; none
of the shadow-rate papers we examined in this line uses a mixture,
particle, or exactly weighted likelihood, and none samples the
posterior with gradient MCMC. That statement must, however, be scoped
to the closure-filter line: outside it, Pericoli and Taboga (2018,
Secs.~2, 4--5) estimate a Black-type max shadow-rate model
(\(r_t=\max\{\bar r_t,\,a+bX_t\}\) with a time-varying bound) by fully
Bayesian data-augmentation MCMC, blockwise random-walk Metropolis
over parameters \emph{and} the entire latent factor path, with no
Gaussian moment-closure filter of any kind, using a neural-network
bond-pricing surrogate trained to sub-basis-point accuracy; their
``nearly exact'' refers to that pricing approximation, not to the
sampler. The exact-likelihood (closure-free) route in this literature
is therefore already occupied. A UKF-plus-HMC route can claim a
gradient-sampler frontier within the closure-filter family, but not a
likelihood-fidelity frontier; the exact constructions below are
measured against Pericoli--Taboga's closure-free standard, not merely
against the filtering papers.

\subsection{A ladder of three targets and the geometry of the
hard-bound variants}
\label{sec:shadow-ladder}

``Exact'' is meaningless without naming what it is relative
to, and this model admits three natural versions that are
\emph{different posteriors}. We name them once and use the names
throughout: \(\mathcal S_1\), the smooth softplus model, i.e.\ the finite
program of \eqref{eq:77}--\eqref{eq:78} with \(K=40\) quadrature of the
softplus map, the model of Section~\ref{sec:shadow-model} exactly as
specified; \(\mathcal C_1\), the finite hard-max model, the same
nodes and weights with \(s_{\ell,\alpha}\) replaced by \(m_\ell\); and
\(\mathcal C_0\), the continuous-maturity hard model, the integral
\(\tau_i^{-1}\int_0^{\tau_i} m_{\ell_c}(f_c(u;x))\,du\) evaluated by
root-bracketed integration to a stated tolerance. Table~\ref{tab:ladder}
summarizes the ladder.

\begin{table}[htbp]\small
\begin{tabularx}{\textwidth}{@{}>{\raggedright\arraybackslash}p{0.16\textwidth}XX@{}}
\toprule
Variant & Definition & In what sense inference under it is exact \\
\midrule
\(\mathcal S_1\) (smooth softplus model) &
The finite program of \eqref{eq:77}--\eqref{eq:78}: \(K=40\)
Gauss--Legendre quadrature of the softplus-bounded forward curve. &
Smooth-model inference: any exact filter or sampler is exact for the
softplus posterior, never for a hard-bound posterior
(Section~\ref{sec:shadow-softplus}). \\
\(\mathcal C_1\) (finite hard-max model) &
The same quadrature nodes and weights with \(s_{\ell,\alpha}\)
replaced by the hard max \(m_\ell\). &
Exact relative to the finite hard-quadrature model: the bound is hard
but maturity averaging is the fixed \(K\)-node rule. \\
\(\mathcal C_0\) (continuous-maturity hard model) &
The continuous integral
\(\tau_i^{-1}\int_0^{\tau_i} m_{\ell_c}(f_c(u;x))\,du\), evaluated by
root-bracketed integration to a stated tolerance. &
Exact relative to the declared integral and tolerance: the bound and
the maturity average are both taken at the level of the continuous
model. \\
\bottomrule
\end{tabularx}
\caption{The ladder of three model variants. Each is a different
posterior; every likelihood, sampler, and result in this case study is
stated relative to one of them.}\label{tab:ladder}
\end{table}

\(\mathcal C_1\) is the natural first hard-bound counterfactual: it
changes exactly one operation in the smooth model. But \(\mathcal C_0\)
and \(\mathcal C_1\) are different posteriors, for a structural reason
rather than a numerical one. Under a simple
transverse root of the bound gap, the moving-boundary terms of the
\(\mathcal C_0\) integral cancel because the two integrand branches
agree at the crossing, so \(y_c^{\mathcal C_0}(\tau_i;x)\) is
continuously differentiable in \(x\) there, whereas the finite
pointwise-max sum of \(\mathcal C_1\) has kinks whenever any node
crosses the bound. Tangencies, endpoint roots, and an identically bound
curve are separate cases requiring their own treatment in the
\(\mathcal C_0\) integration scheme. Root-aware \(\mathcal C_0\)
evaluation and a \(\mathcal C_0\)-versus-\(\mathcal C_1\)
discretization sensitivity experiment are prerequisites before any
event- or kink-aware HMC design should be adopted as the final
hard-ZLB architecture; we specify that experiment here without
running it.

Replace \(s_{\ell,\alpha}\) by \(m_\ell\) in \eqref{eq:78} to obtain
\(\mathcal C_1\). Each yield becomes a nonnegatively weighted sum of
maxima of affine functions of the state,
\begin{equation}\label{eq:81}
 y_c^{\max}(\tau_i;x)=\sum_{k=1}^{K}w_k
 \max\bigl\{\ell_c,\;a_c(s_{ik})^{\mathsf T}x^c\bigr\},
 \qquad s_{ik}=\tau_i v_k,
\end{equation}
which is continuous, convex, and piecewise affine in \(x^c\). The FX
row \eqref{eq:79} is then a difference of such functions plus an affine
basis term: continuous and piecewise affine, though no longer convex.
The complete observation map \(h\) is continuous and piecewise affine
with kinks on the hyperplane arrangement
\(\bigcup_{c,i,k}\{x:a_c(s_{ik})^{\mathsf T}x^c=\ell_c\}\), a
Lebesgue-null set. Because the transition \eqref{eq:75} is a fixed
linear-Gaussian map with no branch, and the observation noise is
full-support Gaussian, the joint posterior over \((\vartheta,x_{0:T})\),
equivalently over \((\vartheta,x_0,\eta_{1:T})\) in the
non-centered chart of \eqref{eq:56}, is continuous, is smooth off a
finite hyperplane arrangement, and has no atoms and no density jumps.
In the classification of Section~\ref{sec:geometries} this is geometry one: a
kink target, not a jump target. Consequently ordinary Metropolized
leapfrog HMC is valid on the \(\mathcal C_1\) target by
Section~\ref{sec:kink-validity}, the refraction machinery of Section~\ref{sec:event}
degenerates (\(\Delta U=0\)), and discontinuous or mixed HMC address
structure that this model simply does not have. (The \(\mathcal C_0\)
target is smoother still at transverse roots, per the ladder discussion
above; the geometry-one classification covers both hard-bound
variants.)

The classification is a statement about the model as specified, and it
is fragile in three specific directions, each of which a modeler may
adopt deliberately. If the factor dynamics switch at the bound (a
binding-dependent \(\Phi\) or \(Q\)), the transition acquires a
genuine branch and the target moves to geometry two, with everything
that Section~\ref{sec:geometries}, Section~\ref{sec:event}, and Section~\ref{sec:target-det}
say about deterministic jumps. If a policy rate is added as an
observation \emph{without} measurement error,
\(i_t=\max\{\ell,f_c(0;x_t)\}\) observed exactly, the observation law
acquires an atom at \(\ell\) and the likelihood becomes the censored
mixture of Section~\ref{sec:ukf-censored} with \(V\to0\). And if the lower
bound \(\ell\) is estimated, as in Lemke and Vladu (2017), the kink
location becomes parameter-dependent, which leaves continuity intact
but makes the kink arrangement move with \(\vartheta\). A modeler must
decide which of these extensions, if any, is in play; the geometry
claim above holds for none of them automatically.

One structural lemma sharpens everything downstream. Fix \(x^c\) and
write the bound gap along the horizon as
\begin{equation}\label{eq:82}
 g(s)=f_c(s;x)-\ell_c=(L-\ell_c)+e^{-\lambda_c s}(S+C\lambda_c s),
 \qquad
 g'(s)=\lambda_c e^{-\lambda_c s}\bigl(C-S-C\lambda_c s\bigr),
\end{equation}
with \((L,S,C)=x^c\). If \(L-\ell_c=0\), \(S=0\), and \(C=0\), the gap
is identically zero and the curve sits exactly on the bound at every
horizon; this degenerate case must be handled by convention (it is a
Lebesgue-null event under the Gaussian predictive, though synthetic
test cases can be constructed to land on it). Away from it, the bracket
in \(g'\) is affine in \(s\), so \(g'\) changes sign at most once;
\(g\) is monotone or single-peaked/troughed on \((0,\infty)\) and
therefore has \emph{at most two zeros}. The binding set
\(\{s:f_c(s;x)\leq\ell_c\}\) along any Nelson--Siegel forward curve is
an interval, the complement of an interval, empty, or everything. Two
consequences follow. The reachable binding patterns over the pooled
quadrature horizons \(\{s_{ik}\}\) are \emph{interval patterns},
\(O(\text{grid}^2)\) many rather than \(2^{\text{grid}}\); and the
per-country kink cells in factor space are polyhedra cut by constraints
whose normals \(a_c(s)\) span at most that country's three factor
dimensions, so per-country cell probabilities reduce to at most
three-dimensional Gaussian integrals. The \emph{joint} two-country cell
probability required by the evidence decomposition below is generically
six-dimensional: the FX row and the measurement update correlate the
two factor blocks, so the joint mass factorizes into two
three-dimensional integrals only if a factorization of the relevant
covariance is proved for the case at hand.

\subsection{An exact cell identity and the bootstrap filter that carries
the likelihood}
\label{sec:shadow-filter}

Because the observation map of \(\mathcal C_1\) is exactly affine on each
polyhedral cell of the kink arrangement, a beautiful idea suggests itself:
the one-step filtering problem should decompose exactly over cells into
conjugate Gaussian updates. This
section develops that identity in full, because it is genuinely exact and
it organizes everything one can know about the model's local structure,
and then explains why three specific obstacles keep it from being
an implementable exact algorithm. The pedagogical arc is deliberate:
first the ideal identity, then precisely where it breaks, then the
robust alternative that actually carries the inferential weight. That
alternative is the plain bootstrap particle filter, and we state it
first so the reader knows where the argument lands.

For \(\mathcal C_1\), the workhorse likelihood estimator is the
\textbf{bootstrap particle filter}: per particle, draw the state from
the exact Gaussian transition and weight by the exact observation
density,
\begin{equation}\label{eq:82a}
 x_t^j\sim N\bigl((I-\Phi)\bar\theta+\Phi x_{t-1}^{A_t^j},\,Q\bigr),
 \qquad
 v_t^j=N\bigl(y_t;\,h_{\max}(x_t^j),\,R\bigr),
\end{equation}
where \(h_{\max}\) stacks the \(\mathcal C_1\) yield rows
\eqref{eq:81} and the FX row \eqref{eq:79}. Both operations are
available in closed form, the estimator \eqref{eq:13} is nonnegative
and unbiased under the standard support and integrability conditions
of Section~\ref{sec:pf-generic}, and PMMH then carries the ordinary
particle-MCMC extended-target guarantee of Section~\ref{sec:pm-target}. This
route is less elegant than the cell decomposition but is implementable
today, unbiased for the named finite model, and its Monte Carlo variance is
directly measurable from likelihood replicates.

Now the ideal. Conditional on a binding pattern \(B\), an assignment
of bind/slack to every quadrature horizon, the observation map is
affine, \(h_B(x)=H_Bx+d_B\), on the polyhedral cell
\begin{equation}\label{eq:83}
 C_B=\bigl\{x:\;a_c(s_{ik})^{\mathsf T}x^c\leq\ell_c\ \text{for}\
 (c,i,k)\in B,\ \ >\ell_c\ \text{otherwise}\bigr\}.
\end{equation}
Let the one-step predictive be Gaussian,
\(x_t\mid\mathcal F\sim N(m,P)\); this is exactly true per particle, since
\eqref{eq:75} gives
\(x_t\mid x_{t-1}\sim N((I-\Phi)\bar\theta+\Phi x_{t-1},Q)\). The exact
one-step posterior and evidence then decompose over cells with no
approximation:
\begin{equation}\label{eq:84}
 p(x_t\mid y_t,\mathcal F)\;\propto\;
 \sum_{B}\mathbf 1\{x_t\in C_B\}\,
 N\!\bigl(y_t;H_Bm+d_B,\,S_B\bigr)\,
 N\!\bigl(x_t;m_B^+,P_B^+\bigr),
\end{equation}
where each cell carries the standard conjugate update
\begin{equation}\label{eq:85}
 S_B=H_BPH_B^{\mathsf T}+R,\quad
 K_B=PH_B^{\mathsf T}S_B^{-1},\quad
 m_B^+=m+K_B(y_t-H_Bm-d_B),\quad
 P_B^+=P-K_BS_BK_B^{\mathsf T},
\end{equation}
in which \(S_B\) is invertible for every cell because the measurement
covariance \(R\) is positive definite and \(H_BPH_B^{\mathsf T}\) is
positive semidefinite, and the exact evidence is the polyhedrally
truncated mixture mass
\begin{equation}\label{eq:86}
 p(y_t\mid\mathcal F)=\sum_B Z_B,
 \qquad
 Z_B=N\!\bigl(y_t;H_Bm+d_B,\,S_B\bigr)\,
 \Pr_{x\sim N(m_B^+,P_B^+)}\!\bigl(x\in C_B\bigr).
\end{equation}
Equations \eqref{eq:84}--\eqref{eq:86} are the hard-bound
observation-map counterpart of the branch-split update
\eqref{eq:36}--\eqref{eq:45}: the branch structure has moved from the
transition maps \(F_r\) into the observation cells \(C_B\), the
truncated moments of \eqref{eq:21}--\eqref{eq:28} reappear whenever a
single constraint dominates, and by the crossing lemma of
Section~\ref{sec:shadow-ladder} the sum over \(B\) runs over interval patterns
only.

Inside a particle filter the identity yields an \emph{ideal} fully
adapted construction. Per particle \(j\) with ancestor \(x_{t-1}^j\),
the predictive is exactly Gaussian, so \eqref{eq:84}--\eqref{eq:86}
hold exactly with
\((m,P)=((I-\Phi)\bar\theta+\Phi x_{t-1}^j,\,Q)\), and the
conditionally optimal proposal is, in the ideal,
\begin{equation}\label{eq:87}
 q^\star(x_t\mid x_{t-1}^j,y_t)=p(x_t\mid x_{t-1}^j,y_t),
 \qquad
 v_t^j=p(y_t\mid x_{t-1}^j)=\sum_B Z_B\bigl(x_{t-1}^j\bigr),
\end{equation}
sampled by drawing a cell with probability \(\propto Z_B\), then
drawing from the cell-restricted Gaussian
\(N(m_B^+,P_B^+)\mathbf 1_{C_B}\). The point-independence of the
incremental weight \eqref{eq:87} is the fully adapted property of
\eqref{eq:12b}--\eqref{eq:13}, and \emph{if} every \(Z_B\) is computed
exactly and the cell-restricted draw is exact, the resulting likelihood
estimator is nonnegative and unbiased for the \(\mathcal C_1\) model by
the argument of Section~\ref{sec:pf-generic}.

So why is this not the algorithm of record? Three specific reasons,
each a numerical obstacle rather than a conceptual flaw. First,
the cell masses \(\Pr_{x\sim N(m_B^+,P_B^+)}(x\in C_B)\) are
multivariate Gaussian polytope probabilities, generically
six-dimensional across the two correlated factor blocks per the
crossing-lemma discussion, and these are numerical computations in
the sense of Genz and Bretz (2009), not closed-form operations.
Second, restricting the sum to patterns with ``non-negligible'' mass is
a \emph{pruning approximation}: dropping any cell changes
\(\sum_B Z_B\), the cell-selection probabilities, and the weight, so
the result is exact only if the omitted masses are certified to be
exactly zero. Third, exact draws from the cell-restricted Gaussian are
themselves nontrivial: rejection sampling is exact in principle but can
be unusably inefficient in the tails; the exact piecewise-quadratic HMC
of Pakman and Paninski (2014) supplies an \emph{invariant kernel} for
the truncated Gaussian, not an independent exact draw after a finite
trajectory, so substituting one trajectory for an exact draw changes
the proposal law and forfeits the constant weight \eqref{eq:87}; and
Botev's (2017) minimax-tilting method is the modern benchmark for both
the probabilities and the draws but carries its own numerical error,
which must be quantified and folded into the proposal before any
exactness language attaches to a cell implementation. A related
limitation applies beyond the particle setting: \eqref{eq:84}--\eqref{eq:86}
do \emph{not} give a closed-form multi-step filter, because propagating
a cell-restricted Gaussian through \eqref{eq:75} leaves the Gaussian
family, so an exact density filter must carry a growing
truncated-mixture representation, and any pruning or moment collapse
\eqref{eq:47}--\eqref{eq:48} is a declared approximation, exactly as in
Section~\ref{sec:ukf-branchwise}.

The resolution is the one already announced: until those numerical
components are controlled, the cell construction is a \emph{proposal
and diagnostic layer} (an implementation that prunes cells or
approximates draws must use its actual proposal density and the
corresponding importance ratio in \eqref{eq:12b}, not the ideal
constant weight \eqref{eq:87}), while the bootstrap route
\eqref{eq:82a} carries the likelihood. This positioning is the Aruoba
et al.\ conditionally-optimal construction transplanted from a
piecewise-linear transition to a hard-bound observation map: simpler
here because the transition never branches, but gated on the same
numerical error control.

The Tobit line of the censored-filtering literature sits on the other
side of a noise-ordering distinction that decides its relevance here.
In the Tobit-type-1 measurement model surveyed
by Geng et al.\ (2021), censoring is applied \emph{after} the noise,
\(y=\max\{c,\,Hx+u\}\), which produces an atom at the censoring limit
and is the setting of the Tobit Kalman filter (Allik et al.\ 2016; we
cite this from metadata, having been unable to obtain the full text).
Our model's yields add noise \emph{after} the max, so they have no
atom, and the mixture structure lives in the state, not in the
observation support. The two orderings coincide only in the exactly
observed policy-rate extension of Section~\ref{sec:shadow-ladder}, which is
where the Tobit machinery would become the right reference.

\subsection{The softplus route: bias bounds and two readings of
\texorpdfstring{$\alpha$}{alpha}}
\label{sec:shadow-softplus}

How far is the smooth model \(\mathcal S_1\) from the hard model
\(\mathcal C_1\)? The pointwise gap between the two scalar maps is
explicit:
\begin{equation}\label{eq:88}
 0\;<\;s_{\ell,\alpha}(u)-m_\ell(u)
 =\alpha\log\!\left(1+e^{-|u-\ell|/\alpha}\right)
 \;\leq\;\alpha\log 2,
\end{equation}
with the maximum attained exactly at the kink and exponential decay
away from it: at \(|u-\ell|\geq5\alpha\) the gap is below
\(6.8\times10^{-3}\,\alpha\). Averaging \eqref{eq:88} through the
quadrature \eqref{eq:78} bounds the yield gap uniformly,
\begin{equation}\label{eq:89}
 0\;<\;y_c^{\alpha}(\tau_i;x)-y_c^{\max}(\tau_i;x)\;\leq\;\alpha_c\log 2,
\end{equation}
and the bound is attained only when the whole forward path sits at the
kink; in general the gap is \(\alpha\log2\) times the quadrature mass
within \(O(\alpha)\) of the bound plus an exponentially small
remainder. With the working example's constants the worst-case yield
gaps are \(1.04\times10^{-3}\) (domestic) and \(0.69\times10^{-3}\)
(foreign) in annualized rate units, roughly ten and seven basis
points. The gap is moreover \emph{one-signed}: softplus
yields always exceed hard-censoring yields, so fitted factors and
long-run means absorb a systematic downward shift near binding rather
than symmetric noise. The per-observation log-likelihood perturbation
obeys
\begin{equation}\label{eq:90}
 \bigl|\log N(y;h_\alpha,R)-\log N(y;h_{\max},R)\bigr|
 \;\leq\;
 \bigl\|R^{-1/2}(y-h_{\max})\bigr\|\,
 \bigl\|R^{-1/2}\Delta\bigr\|
 +\tfrac12\bigl\|R^{-1/2}\Delta\bigr\|^2,
\end{equation}
with \(0\le\Delta_i\le\alpha\log2\) componentwise, which is small per
observation. The signed difference itself is, for
\(\Delta=h_\alpha-h_{\max}\),
\begin{equation}\label{eq:90a}
 \log N(y;h_\alpha,R)-\log N(y;h_{\max},R)
 =(y-h_{\max})^{\mathsf T}R^{-1}\Delta
 -\tfrac12\,\Delta^{\mathsf T}R^{-1}\Delta,
\end{equation}
and this can take \emph{either sign} even though the map shift
\(\Delta\) is one-sided on the yield rows: the sign depends on the
residual \(y-h_{\max}\), and the FX row's \(\Delta\) component is a
domestic-minus-foreign combination that can itself be signed either
way. In summary: the \emph{mean-map}
discrepancy is one-sided and may induce a systematic fitted-factor
adjustment near binding, while the likelihood and posterior
perturbations take either sign and should be reported as computed signed
quantities under
\eqref{eq:90}--\eqref{eq:90a}.
With measurement-error standard deviations of a few basis points,
a ten-basis-point systematic map gap is not negligible relative to the
noise scale, and \eqref{eq:90}--\eqref{eq:90a} quantify it in exactly the
units, log-likelihood, in which it affects inference.

One might hope to shrink the gap by sharpening \(\alpha\), but this is
not free for gradient samplers. The scalar curvature is
\begin{equation}\label{eq:91}
 s_{\ell,\alpha}''(u)=\frac{\sigma(z)\bigl(1-\sigma(z)\bigr)}{\alpha}
 \leq\frac1{4\alpha},
 \qquad z=\frac{u-\ell}{\alpha},
\end{equation}
so the observation Hessian acquires ridges of height \(O(1/\alpha)\)
concentrated in an \(O(\alpha)\)-neighborhood of the kink, and leapfrog
stability suggests a local step size \(\epsilon=O(\sqrt\alpha)\) there.
That scaling is a local stability heuristic under a fixed metric and
comparable posterior curvature, not a universal necessity theorem: the
binding step size depends on the chart, the mass matrix, the
trajectory, and how much time the sampler spends near the bound. The
qualitative conclusion survives the qualification: the \(\alpha\to0\)
limit of the smooth route is \emph{harder} for HMC than the
\(\mathcal C_1\) hard-bound model itself, whose one-sided curvature is
bounded and whose kink cost is the \(O(\epsilon)\) energy error of
\eqref{eq:56a}. Sharpening the softplus to approximate the hard bound
is the wrong direction; the right choices are either a genuinely smooth
model or the exact machinery of Section~\ref{sec:shadow-filter}.

That leaves two defensible readings of \(\alpha\), and they support
different claims. Read as a \emph{numerical approximation} to the hard-ZLB model,
\(\alpha\) needs an error protocol: report
\eqref{eq:89}--\eqref{eq:90} alongside binding fractions, rerun matched
synthetic-data estimations at \(\alpha/2\) and \(2\alpha\) (in our
experiments this is implemented by re-running matched synthetic-data
estimations at halved and doubled \(\alpha\)), and require the
posterior functionals of interest to move by less than their Monte
Carlo uncertainty before any hard-ZLB language is used. Read as a
\emph{structural smooth-transition model} in the sense argued
empirically by Opschoor and van der Wel (2024), whose evidence is that
yields leave the bound gradually and the smooth map may therefore be the
better description,
\(\alpha\) is a model constant or even an estimable parameter, no bias
language applies, and every claim must name the smooth model as the
target. The hybrid to avoid, under both readings, is sampling the smooth
model while describing the result as inference under the ZLB: the two
posteriors differ by exactly the gaps computed above.

\subsection{Identification at the bound}
\label{sec:shadow-ident}

Complicating the likelihood's geometry is not the only thing the bound
does; it also erodes what the data can say. Differentiating \eqref{eq:78} gives the
yield sensitivity to the factors,
\begin{equation}\label{eq:92}
 \frac{\partial y_c(\tau_i;x)}{\partial x^c}
 =\sum_{k=1}^{K}w_k\,
 \sigma\!\left(\frac{f_c(s_{ik};x)-\ell_c}{\alpha_c}\right)
 a_c(s_{ik}),
\end{equation}
with the logistic factors replaced by indicators
\(\mathbf 1\{f>\ell\}\) in the \(\mathcal C_1\) hard-bound model. When
the forward curve sits far below the bound over the horizons that
matter for maturity \(\tau_i\), the corresponding rows of the
observation Jacobian \(J\) vanish, and the per-period observed Fisher
information \(J^{\mathsf T}R^{-1}J\) about the shadow factors
degenerates in exactly the binding directions. During a long binding
spell the shadow-factor posterior is driven by the prior and the
transition \eqref{eq:75} alone, and the long-run means \(\bar\theta\)
of the binding country are informed only through occasional excursions
of the forward curve above the bound. This is the formal counterpart of
two empirical findings: Bauer and Rudebusch (2016) document that
shadow-rate estimates are highly sensitive to the assumed lower-bound
value and to model specification, and Christensen and Rudebusch (2015)
report the same sensitivity for shadow-rate paths. The practical
consequences are direct. Priors on \(\bar\theta\) do real work during
binding spells, so reporting them is part of reporting the result. And the
identification degeneracy makes the sampler ill-conditioned in exactly the same
directions: the HMC mass matrix appropriate to binding-era geometry is
very different from the one appropriate off the bound, so step-size and
mass-matrix tuning performed on off-bound data, or on another
model, horizon, or particle count, will not transfer to binding
regimes and must be redone for the configuration actually being run.
Any empirical claim from this model therefore comes with identification
diagnostics along the binding directions, since those directions are where
the data stop speaking and the prior takes over.

\subsection{Five routes and what each delivers}
\label{sec:shadow-routes}

We can now assemble the case study's conclusions as a menu of five
inference routes. They differ in which target on the ladder of
Table~\ref{tab:ladder} they address and in the strength of the
guarantee they carry, not merely in cost; Table~\ref{tab:routes}
summarizes.

\begin{table}[htbp]\small
\begin{tabularx}{\textwidth}{@{}>{\raggedright\arraybackslash}p{0.16\textwidth}XX@{}}
\toprule
Route & Target & What inference under this route delivers \\
\midrule
(i) SR-UKF marginal likelihood + HMC/NeuTra &
The UKF likelihood as a stated approximate functional of
\(\mathcal S_1\) &
Exact MCMC for its own target, the UKF likelihood rather than the model
likelihood; this mirrors the estimator of Kim and Priebsch (2013). Its
single-Gaussian closure discards branch mass exactly as in
Section~\ref{sec:ukf-branchloss}; its distance from route (iii) on matched
binding scenarios is a required diagnostic, not an assumption. \\
(ii) Softplus model with any exact filter/sampler &
\(\mathcal S_1\) &
Legitimate as a structural model (the Opschoor--van der Wel reading)
or as an approximation under the \eqref{eq:89}--\eqref{eq:90a}
protocol; never a hard-ZLB claim by itself. \\
(iii) Bootstrap hard-bound particle filter \eqref{eq:82a} + PMMH &
\(\mathcal C_1\) &
The likelihood workhorse: exact Gaussian transition draws, exact
observation-density weights, a nonnegative unbiased estimator, and
exact extended-target MCMC by Section~\ref{sec:pm-target}. Delivers valid
inference for the finite hard-quadrature posterior. Supports PMMH, not
pseudo-marginal HMC, absent the differentiability conditions of
Section~\ref{sec:pmhmc}; correlated pseudo-marginal (Deligiannidis, Doucet,
and Pitt 2018) is the variance/mixing comparator, not a
differentiability fix. \\
(iv) Joint \((\vartheta,x_0,\eta_{1:T})\) HMC on \(\mathcal C_1\) &
The \(\mathcal C_1\) model, sampled in the non-centered chart of
\eqref{eq:56} &
Valid by Section~\ref{sec:kink-validity} because the target is a continuous
piecewise-quadratic-plus-smooth kink target; needs no filter at all.
Gradient cost is one simulation sweep, and the linear-Gaussian prior
structure supplies the natural preconditioner. Kink crossings cost
acceptance per \eqref{eq:56a}, not correctness. \\
(v) Cell-adapted fully adapted proposal
\eqref{eq:84}--\eqref{eq:87} inside a particle filter &
\(\mathcal C_1\) &
Conditional: delivers exact fully adapted filtering only once the
Genz--Bretz polytope-probability and Botev truncated-draw components
are benchmarked and carried with quantified numerical error; with
pruning or approximate draws it must use its actual proposal density
and importance ratio, and until then it is a proposal/diagnostic layer
for route (iii). \\
\bottomrule
\end{tabularx}
\caption{Five inference routes for the two-country shadow-rate model,
the ladder target each addresses, and what inference under each
delivers.}\label{tab:routes}
\end{table}

The branch-split mixture UKF of Section~\ref{sec:ukf} remains what it was
there: a fast diagnostic and a proposal ingredient, with the likelihood of
record carried by route (iii). Routes (iii) and (iv) target the same
\(\mathcal C_1\) model and must agree within Monte Carlo error on
synthetic data; that cross-check is the cheapest end-to-end
correctness test available for this model, and it is executable at
small parameter-recovery scale before any full-data run. The
continuous-maturity target \(\mathcal C_0\) enters through the
discretization sensitivity experiment of Section~\ref{sec:shadow-ladder}
before any event- or kink-aware design is adopted.

\section{Case study II: a New Keynesian model with an occasionally binding policy rate}
\label{sec:nk}

Our second case study turns from term structures to the model class that
motivates most of the ZLB estimation literature: a small New Keynesian economy
whose policy rate is truncated at a lower bound and whose equilibrium is
computed by a piecewise-linear method. Likelihoods built on such solutions are
folklore-famous for being ``spiky'': practitioners report jagged likelihood
surfaces, optimizers that stall on cliffs, and gradient-based samplers that
reject in bursts. It is tempting to blame the $\max$ operator in the policy
rule. The analysis of this section shows that the blame is misplaced. With a
unique verified solution, the piecewise-linear ZLB model is a \emph{kink}
target in states and shocks: continuous, piecewise smooth, and well within
the reach of the methods of Section~\ref{sec:hmc}. The genuine discontinuities
practitioners meet come from the solution \emph{correspondence},
nonexistence and multiplicity of verified regime paths, and from the
selection rule imposed on it. We derive this claim in full on the smallest
model that exhibits it, then draw out what it implies for how an estimation
exercise must be posed.

\subsection{Solution operators for the piecewise-linear class}
\label{sec:nk-solvers}

Take the benchmark model class: the three-equation New Keynesian model with a
Taylor rule truncated at the bound, written in the notation of
\eqref{eq:1}--\eqref{eq:3} and solved by a piecewise-linear method. Three
solution operators dominate the literature. OccBin (Guerrieri and Iacoviello
2015, Sec.~2) guesses a regime path, solves the two linear systems backward
from a terminal reference-regime date, simulates forward, and iterates until
the implied inequalities verify the guess; the conditional filtering
counterpart of this construction is \eqref{eq:6}--\eqref{eq:11}. Holden (2016)
reformulates the bound as anticipated news shocks added to the unconstrained
linear solution and implements the search in DynareOBC; the anticipated-shock
route originates with Holden and Paetz (2012), whom we cite for orientation.
Boehl (2022) derives a closed-form expression for the whole trajectory given
the guessed spell durations, which reduces the iteration to a fast search over
two integers. All three produce, when they converge to a verified path, the
same object: a solution that is linear on each regime cell of the state-shock
space.

Holden (2023, Secs.~2--3) supplies the algebra that makes the solution
correspondence analyzable; pause on it, because it converts a
vague worry about ``convergence failures'' into a precise mathematical
question. Stack the bound violations of the unconstrained solution into
$q\in\mathbb R^{T}$ and let $M\in\mathbb R^{T\times T}$ collect the responses
of the bounded variable to unit news shocks; imposing the bound over a horizon
of $T$ periods is the linear complementarity problem
\begin{equation}\label{eq:93}
 y\geq0,\qquad q+My\geq0,\qquad y^{\mathsf T}(q+My)=0,
\end{equation}
whose solution $y$ gives the news-shock loadings that hold the variable at the
bound during the binding spell. Existence and uniqueness of the verified path
are exactly existence and uniqueness for \eqref{eq:93}; Holden's central
result is that the solution is unique for \emph{every} $q$ if and only if $M$
is a P-matrix (all principal minors positive), and that standard New Keynesian
models fail this condition, having multiple perfect-foresight paths that
escape the bound.

The P-matrix property is a uniqueness result at each parameter point, and
continuity of the solution map is a separate question about how the solution
moves \emph{between} points. Pointwise uniqueness of the LCP solution does
not by itself establish that the solution map
$\vartheta\mapsto y(q(\vartheta),M(\vartheta))$ is continuous. The usable
regularity statement is parametric: for a \emph{fixed finite horizon} $T$ and
terminal condition, with $q(\vartheta)$ and $M(\vartheta)$ continuous on a
stated parameter domain over which $M(\vartheta)$ remains a P-matrix, the
unique LCP solution is continuous in $\vartheta$ by standard parametric-LCP
regularity. Any of a changing horizon, a changed terminal condition, or loss
of the P-matrix property on the domain voids the conclusion and can create
genuine discontinuity even under pointwise uniqueness; an estimation exercise
should therefore state which of these conditions has been verified, and on
which domain.

\subsection{A fully derived multiplicity example and the discontinuity mechanism}
\label{sec:nk-mult}

Holden's (2023, Sec.~2.2) lagged-response economy is the smallest model that
shows the mechanism, and we reproduce it in full because the rest of the
case study leans on it. The Fisher equation with a constant real rate $r>0$
and a truncated Taylor rule with lagged response give
\begin{equation}\label{eq:94}
 r+\pi_{t+1}=i_t=\max\{0,\;r+\phi\pi_t-\psi\pi_{t-1}\},
 \qquad \phi-\psi>1,\ \psi\in(0,1),
\end{equation}
with $\pi_0$ given. Away from the bound the stable solution is
$\pi_t=A\pi_{t-1}$ with $A^{2}-\phi A+\psi=0$; taking $\phi=2$ for
transparency, $A=1-\sqrt{1-\psi}\in(0,1)$ and the useful identity
$A(\phi-A)=\psi$ holds. Holden shows the economy can be at the bound for at
most the first period on any path that escapes, so $\pi_2=A\pi_1$, and
substituting into \eqref{eq:94} at $t=1$ leaves a scalar instance of
\eqref{eq:93}:
\begin{equation}\label{eq:95}
 0=\max\bigl\{-r-A\pi_1,\;(\phi-A)\pi_1-\psi\pi_0\bigr\}.
\end{equation}
Both branches can be solved explicitly. The slack branch gives the fundamental
solution $\pi_1^{\mathrm f}=\psi\pi_0/(\phi-A)$, admissible when
$r+A\pi_1^{\mathrm f}\geq0$, that is when $\pi_0\geq-r/A^{2}$ by the identity
above. The binding branch gives $\pi_1^{\mathrm b}=-r/A$, admissible when
$(\phi-A)\pi_1^{\mathrm b}\leq\psi\pi_0$, which is the \emph{same} condition
$\pi_0\geq-r/A^{2}$. Therefore
\begin{equation}\label{eq:96}
\begin{aligned}
 \pi_0>-\frac{r}{A^{2}}&:\ \text{two solutions},\\
 \pi_0=-\frac{r}{A^{2}}&:\ \text{one},\\
 \pi_0<-\frac{r}{A^{2}}&:\ \text{none returning to the standard steady state}.
\end{aligned}
\end{equation}

Multiplicity here covers an open half-line of
initial states rather than a knife-edge, and Holden proves it becomes pervasive in richer New Keynesian
models. Now observe where the discontinuity actually sits. Each branch map is
\emph{continuous} in $(\pi_0,r,\psi)$ ($\pi_1^{\mathrm f}$ is linear in
$\pi_0$, $\pi_1^{\mathrm b}$ is constant), so the discontinuity is created
entirely by \emph{selection}: any deterministic rule $\sigma$ that picks a
branch defines a transition kernel
\begin{equation}\label{eq:97}
 \pi_1=T_{\sigma(\pi_0,\vartheta)}(\pi_0,\vartheta),
 \qquad
 T_{\mathrm f}\neq T_{\mathrm b}\ \text{on the whole multiplicity region},
\end{equation}
which jumps by an $O(1)$ amount wherever $\sigma$ switches branch. A
likelihood built on \eqref{eq:97} inherits that jump in the state, and,
because the multiplicity boundary $-r/A(\psi)^{2}$ and any solver-dependent
selection move with $\vartheta$, also in the parameters. This is the formal
content of the folklore that piecewise-linear ZLB likelihoods ``spike'': the
spike is not caused by the $\max$, whose solution branches are individually
continuous, but by an undeclared, possibly solver-order-dependent selection
rule on a positive-measure multiplicity region, plus outright nonexistence
below the threshold. In the taxonomy of Section~\ref{sec:geometries} the object is
geometry four until a selection law is declared; after a deterministic
declaration it is geometry two along the selection switch set; and only the
stochastic completion of the next subsection makes it geometry three, where
the mixed-support kernels of Section~\ref{sec:mixed} and Section~\ref{sec:pgas} legitimately
apply.

\subsection{The stochastic completion and the sunspot literature}
\label{sec:nk-sunspot}

What is the statistically coherent repair? Make selection part of the model.
Attach to every multiplicity region a probability law over its branches
(in \eqref{eq:95}, a probability $p_t$ of the bound branch, possibly Markov in
a latent sunspot state $s_t$), and the target becomes the mixed-support
posterior \eqref{eq:73} with $p(s_t\mid s_{t-1},x_{t-1},\vartheta)>0$ on both
branches. The construction machinery for linear models under indeterminacy is
well established: Lubik and Schorfheide (2003) show that under indeterminacy
the endogenous forecast errors are not uniquely determined by fundamentals,
and sunspot shocks enter as additional structural disturbances with estimable
loadings. The published instance for the ZLB itself is Aruoba, Cuba-Borda, and
Schorfheide (2018), which we cite for orientation, having been unable to
obtain the full text; they build a New Keynesian model whose
targeted-inflation and deflation regimes are selected by a Markov sunspot and
estimate it with a particle filter; Holden (2023) explicitly contrasts his
within-steady-state multiplicity with that steady-state-switching literature.
The consequence for inference is structural: the sunspot completion is the
\emph{only} reading under which regime-path Gibbs, PGAS, or mixed HMC are
valid, because only there does the regime have positive conditional
probability on both values; under a deterministic selection they collapse to
the zero-support situation of \eqref{eq:5a}.

\subsection{Estimation baselines}
\label{sec:nk-baselines}

Three filter families define current estimation practice for this class, and
all are approximations relative to the particle authority of
Section~\ref{sec:particle}. The piecewise Kalman filter of Section~\ref{sec:linear} conditions
on the verified regime path. Boehl and Strobel (2023) estimate medium-scale
ELB models with an ensemble Kalman filter (ensemble members updated by
linear shifting rather than reweighting) on top of the Boehl (2022) solver,
and compare against the inversion filter; like every Gaussian-closure method,
the EnKF inherits the branch-mass deficiency of Section~\ref{sec:ukf-branchloss} at
the bound, now with sampling noise on top. Holden (2017) proposes an extended
skew-t cubature Kalman filter with dynamic state-space reduction, tracking
third and fourth moments precisely because the Gaussian closure loses the
censoring asymmetry. A serious estimation exercise in this class should
therefore report the piecewise Kalman filter, an ensemble Kalman filter arm,
an inversion-filter arm (Cuba-Borda, Guerrieri, Iacoviello, and Zhong 2019,
cited here for orientation), and the bootstrap particle filter on the same
data sets, with the particle route as the authority, mirroring the route
comparison of Section~\ref{sec:shadow-routes} in the first case study.

\subsection{Designing an estimation exercise}
\label{sec:nk-design}

For this model class, then, the posterior is not
even well-defined until a number of modeling choices, usually left implicit
in software defaults, have been made explicit. We spell them
out, in the notation of \eqref{eq:1}--\eqref{eq:3}, because each one changes
the target rather than merely the algorithm.

First, the bound and the truncated policy rule themselves: which variable is
bounded, at what level, and how the notional rule behaves while the bound
binds (in particular, whether it responds to lagged states, which is exactly
what generated multiplicity in Section~\ref{sec:nk-mult}). Second, the expectation
formation and the terminal condition: perfect foresight over what horizon,
returning to which reference regime. The parametric-LCP regularity result of
Section~\ref{sec:nk-solvers} holds for a fixed horizon and terminal condition,
so changing either mid-estimation changes the continuity properties
of the likelihood along with it. Third, the solution operator and its verification loop:
whichever of the OccBin, news-shock, or closed-form iterations is used, the
estimated object is the \emph{verified} path, and an unverified guess accepted
by a lenient stopping rule is a different, and discontinuously different,
object.

Fourth, the uniqueness domain: the LCP matrix $M$ of \eqref{eq:93} for the
chosen horizon, with a computed P-matrix verdict on the parameter region the
posterior will actually explore, so that the domain on which the solution is
unique is known rather than hoped. Fifth, a selection law on every
multiplicity region (deterministic and stated, or stochastic with its
sunspot process), since \eqref{eq:97} shows that the likelihood jumps
wherever an unstated selection switches branch. Sixth, the value the
likelihood assigns on nonexistence regions, where no verified path returns to
the steady state: zero, a penalty, or an enlarged model in which such
histories have positive probability; each choice produces a different
posterior and a different support. Seventh, the measurement equation and error
structure, which determine whether the state-level kinks and jumps analyzed
above survive into the observed-data likelihood at all.

Finally, the design should include a deliberately tiny version of the model,
small enough that the regime path can be enumerated exhaustively,
on which the piecewise Kalman filter, the ensemble arm, and the particle
authority can be compared against exact enumeration. Once these choices are
frozen, the geometry classification of Section~\ref{sec:geometries} can be executed on
the model actually being estimated rather than on its informal description,
and the guidance of Section~\ref{sec:guidance} applies.

\section{Solver-induced and code-induced discontinuities}
\label{sec:solver}

Discontinuity in the \emph{model} is only half the story. A second,
independent family lives in the \emph{evaluation} of the model, and it matters
for gradient samplers even when the intended target is smooth. The distinction
is worth keeping sharp: a model-level jump changes the target and belongs to
Section~\ref{sec:geometries} and Section~\ref{sec:targets}; an evaluation-level branch leaves
the intended target smooth but makes the \emph{computed} log density or its
gradient discontinuous.

Three mechanisms recur in practice. First, matrix-function kernels with
integer branch decisions: automatic differentiation of a matrix exponential
evaluated by Pad\'e approximation with scaling-and-squaring inherits the
integer scaling count, which jumps at norm thresholds; in one production
implementation we examined, the derivative is computed with the scaling
integer frozen, precisely because no smooth-derivative claim holds across the
branch. At such a branch the \emph{value} is continuous up to the
approximation tolerance while the automatic derivative can jump by a finite
amount. Second, factorization conventions: QR and SVD sign and ordering
choices, and square-root filter downdates near rank changes, are
branch-discontinuous functions of their inputs even though the objects they
represent vary smoothly. Third, almost-everywhere gradients: reverse-mode
differentiation of $\max$, elementwise selection, absolute values, and sorting
picks one subgradient at ties, which is harmless for validity by
Section~\ref{sec:kink-validity} \emph{provided the value is continuous}.

The HMC consequences separate cleanly along that proviso. If the computed log
density is continuous and only its computed gradient jumps at evaluation
branches, then the situation is exactly the kink case: the Metropolis
correction keeps the intended target invariant, and the cost is the
$O(\epsilon\|\Delta\nabla U\|)$ acceptance penalty of \eqref{eq:56a} at branch
crossings. If the computed \emph{value} itself jumps (a solver switching
branches beyond its tolerance, an iteration truncated by a state-dependent
stopping rule, a fixed-point loop finding a different verified path), then
the sampler is targeting a discontinuous computed density it does not know
about: rejections concentrate on branch boundaries, apparent reversibility
holds only up to the jump size, and no diagnostic inside the sampler
separates this from a genuine model jump. The requirement for serious runs is
therefore stated at the evaluation level:
\begin{equation}\label{eq:98}
 \sup_{\text{branch boundaries}}
 \bigl|\log\widehat p_{\,\text{branch}\,1}-\log\widehat p_{\,\text{branch}\,2}\bigr|
 \;\ll\;1
\end{equation}
in Metropolis energy units, verified by test, together with branch-stable
derivative implementations (the fixed-scaling matrix-exponential treatment
above is the template), deterministic tie-breaking, and value-continuity tests
across every known integer branch: scaling orders, factorization signs,
regime-path verification loops, and truncation horizons. These are testing
obligations on the likelihood code, not sampler tuning knobs, and the guidance
of the next section assumes they pass.

\section{Practical guidance}
\label{sec:guidance}

We close the technical development with advice to a reader planning Bayesian
inference for a ZLB model. The ordering below matters because each
stage exists to make a later stage interpretable, and each stage leans on
the ladder of standard test models assembled in
Table~\ref{tab:testmodels}.

\subsection*{Define the target and its geometry first}

Before any sampler code, write the model in the notation of
\eqref{eq:1}--\eqref{eq:3}, identify whether the object is a kink, a density
jump, or an explicit discrete path in the taxonomy of Section~\ref{sec:geometries},
and specify how multiple equilibrium solutions are treated.
Section~\ref{sec:shadow} carries this out for the shadow-rate term-structure case
study; Section~\ref{sec:nk-design} lists what the exercise must pin down for the New
Keynesian case, including a computed LCP uniqueness verdict and a stated
selection law. Nothing downstream can compensate for skipping this step: a
sampler cannot be validated against a target that has not been defined.

\subsection*{A ladder of standard test models}
\label{sec:testmodels}

Every stage below presupposes a reference problem whose answer is known by
means independent of the code being tested. The literature has converged on
a reasonably stable set of such problems, and it is worth assembling them in
one place, organized by where the independent answer comes from: a closed
form, exact enumeration, a published reference implementation, or agreement
between methods that are exact for different reasons. Table~\ref{tab:testmodels}
summarizes the ladder; the paragraphs that follow say what each rung
isolates.

\begin{table}[htbp]\small
\begin{tabularx}{\textwidth}{@{}>{\raggedright\arraybackslash}p{0.16\textwidth}>{\raggedright\arraybackslash}p{0.34\textwidth}X@{}}
\toprule
\textbf{Tier} & \textbf{Model} & \textbf{Independent answer and what it isolates} \\
\midrule
0 & Linear-Gaussian state-space model & Kalman filter likelihood in closed form; every particle, sigma-point, and joint-HMC route must reproduce it before a bound is introduced \\
\addlinespace
0 & Conjugate parameter posterior (known states) & Analytic posterior; calibrates the parameter sampler in isolation \\
\addlinespace
1 & Scalar censored-observation model, \eqref{eq:29}--\eqref{eq:35} & Closed-form mixture likelihood; atom versus smoothed-bound handling, branch masses \\
\addlinespace
1 & Affine-boundary truncated Gaussian & Closed-form moments \eqref{eq:21}--\eqref{eq:28}; unit tests for every truncation step \\
\addlinespace
1 & Truncated multivariate Gaussian sampling & Pakman and Paninski (2014) exact HMC and Botev (2017) tilting as reference draws; cell-restricted proposals \\
\addlinespace
2 & Two-state switching linear-Gaussian model, horizon $T\le 12$ & Exact likelihood and smoothing by enumerating all $2^{T}$ regime paths; PF unbiasedness, PGAS invariance, mixed-HMC path moves \\
\addlinespace
2 & Univariate nonlinear growth model (Gordon, Salmond, and Smith 1993) & Dense-grid numerical filter; the canonical particle-filter stress case: severe nonlinearity and multimodality \\
\addlinespace
3 & One-factor shadow-rate model (Gorovoi and Linetsky 2004) & Closed-form bond prices under a hard bound; analytic anchor for the case study of Section~\ref{sec:shadow} \\
\addlinespace
3 & One-node quadrature reduction of $\mathcal C_1$ & Collapses to the Tier-1 censored-scalar likelihood; exact posterior by low-dimensional grid integration \\
\addlinespace
3 & Lagged-rule economy, \eqref{eq:94}--\eqref{eq:96} & Analytic multiplicity boundary $-r/A^{2}$; detection of nonexistence and selection-rule discontinuities \\
\addlinespace
3 & Published OccBin examples (Guerrieri and Iacoviello 2015) & Dynare reference output; regime-path and piecewise-Kalman bookkeeping \\
\addlinespace
3 & Small-scale NK ZLB model (Aruoba et al.\ 2021) & Published COPF likelihood comparisons; conditionally optimal proposal correctness \\
\addlinespace
4 & Joint-distribution tests (Geweke 2004) & Compares two simulators of the same joint law; needs no analytic posterior \\
\addlinespace
4 & Simulation-based calibration (Talts et al.\ 2018) & Rank-uniformity of the true parameter under the fitted posterior across simulated data sets \\
\bottomrule
\end{tabularx}
\caption{A ladder of standard test models for verifying ZLB inference
algorithms, ordered by the source of the independent answer.}
\label{tab:testmodels}
\end{table}

The tiers are ordered so that each new difficulty enters alone. Tier~0
contains no bound at all: on a linear-Gaussian model the Kalman filter gives
the exact marginal likelihood, so a bootstrap filter must match it in
expectation at every particle count, a sigma-point filter must match it to
rounding, and joint HMC on the non-centered chart of \eqref{eq:56} must
reproduce the analytic posterior. A route that fails here has a bookkeeping
error that no ZLB machinery will fix, which is why this rung comes first
even though it exercises nothing specific to the bound.

Tier~1 introduces the bound in the smallest closed-form settings. The
scalar censored-observation model derived in \eqref{eq:29}--\eqref{eq:35} has
a two-term mixture likelihood in closed form, so it separates cleanly the
questions of atom handling, branch-mass computation, and the bias introduced
by smoothing the bound: evaluating the softplus variant on the same data
measures the likelihood gap of \eqref{eq:90}--\eqref{eq:90a} directly against
truth. The affine-boundary truncated moments \eqref{eq:21}--\eqref{eq:28}
serve as unit tests for every truncation step inside a filter, and exact
truncated-Gaussian samplers give reference draws against which any
cell-restricted proposal can be tested for distributional correctness rather
than merely for support.

Tier~2 replaces closed forms with exact enumeration. In a two-state
switching linear-Gaussian model with a short horizon, conditioning on each
of the $2^{T}$ regime paths reduces the model to
\eqref{eq:6}--\eqref{eq:11}, so the exact likelihood and the exact smoothing
distribution are finite sums. This is the natural place to verify that the
particle estimator \eqref{eq:13} is unbiased in repeated sampling, that PGAS
leaves the path posterior invariant, that mixed-HMC path moves accept at the
rates the enumeration predicts, and that a deterministic-regime variant of
the same model rejects regime-only proposals exactly as the support
calculation of Section~\ref{sec:targets} requires. The univariate nonlinear
growth model of Gordon, Salmond, and Smith (1993), with a dense-grid
numerical filter as reference, then adds severe nonlinearity and
multimodality; it contains no bound, but it is the standard stress case for
particle-filter variance and resampling behavior, and a filter that degrades
unexpectedly there will degrade near a binding spell too.

Tier~3 is specific to the lower bound. The one-factor shadow-rate model of
Gorovoi and Linetsky (2004) prices bonds in closed form under a hard bound,
which makes it the natural analytic anchor for the term-structure case
study: the quadrature, censoring, and filtering layers of
Section~\ref{sec:shadow-model} can each be checked against prices that
involve none of them. In the other direction, reducing the $\mathcal C_1$
working example to a single quadrature node makes each yield a censored
affine function of the factors, so the Tier-1 machinery applies verbatim and
the exact posterior is available by grid integration over the low-dimensional
factor space; the bootstrap-filter and joint-HMC routes of
Section~\ref{sec:shadow-routes} must agree with that grid before they are
trusted at $K=40$. On the New Keynesian side, the lagged-rule economy of
\eqref{eq:94}--\eqref{eq:96} has an analytic multiplicity boundary, so it
tests something no smooth benchmark can: whether the estimation code
detects nonexistence, reports multiplicity, and implements the declared
selection rule rather than an accidental solver-order-dependent one. The
published OccBin examples and the small-scale New Keynesian model of Aruoba
et al.\ (2021) close the tier with reference implementations for the
piecewise-linear path iteration and the conditionally optimal proposal.

Tier~4 tests the sampler as a whole and needs no analytic posterior.
Geweke's (2004) joint-distribution test compares two simulators of the same
joint law of parameters and data, one by construction and one through the
posterior sampler, and detects errors anywhere in the chain; simulation-based
calibration (Talts et al.\ 2018) checks that the rank of the true parameter
within the fitted posterior is uniform across data sets simulated from the
prior. Both apply unchanged to ZLB models, and both are particularly
valuable here because they exercise the binding and nonbinding regions in
prior proportion, including the mixed and boundary cases a hand-picked test
set tends to miss. A practical suite runs the ladder in order, promotes a
route only when every applicable rung passes, and keeps the rungs as
regression tests thereafter, since the solver-branch failures of
Section~\ref{sec:solver} have a way of reappearing after unrelated changes.

\subsection*{Validate a linear benchmark}

Reproduce a published piecewise-linear regime path and the corresponding
piecewise-Kalman likelihood, using Dynare/OccBin output as the natural
reference, and compare
every innovation, covariance, regime inequality, and log-likelihood
contribution on synthetic data. The point is not that the linear method is
the final answer (Section~\ref{sec:linear-limits} explains why it cannot be) but that
agreement on the same object rules out an entire class of bookkeeping errors
before nonlinear machinery obscures them. If the verification fixed point is
not unique on the test data, record it: it is direct evidence about the
model's solution correspondence, the object Section~\ref{sec:nk-mult} shows
drives the estimation difficulties.

\subsection*{Build a particle likelihood authority}

Implement the bootstrap particle filter of Section~\ref{sec:pf-generic}, retaining
particle ancestry and likelihood replicates, together with PGAS and the
conditional-optimal special case of Section~\ref{sec:copf} where it applies. Check the
likelihood estimator against exact enumeration in a deliberately tiny model,
then measure its variance across particle counts. For a hard-bound
observation-map model such as the $\mathcal{C}_1$ target of Section~\ref{sec:shadow},
the bootstrap hard-bound filter is the natural authority; adapted-proposal and
moment-closure constructions, the cell decomposition and the branch-split
mixture UKF of Section~\ref{sec:ukf}, are proposal and diagnostic layers to be unit
tested against the affine closed forms of Section~\ref{sec:linear-kalman}, not
replacement authorities. Sensitivity of the posterior to the choice between
nearby model variants (for instance the continuous-maturity versus discrete
hard-max variants of Section~\ref{sec:shadow-ladder}) belongs here, before any
kink-aware sampler design is selected, because it determines whether that
design effort is even aimed at a stable object. A differentiable
particle-filter route should be treated as diagnostic until its extended
target and Hamiltonian correction have been derived, per Section~\ref{sec:dpf}.

\subsection*{Only then, target-specific exact kernels}

Which exact kernel to build depends on the support structure established in the first
stage, and evidence does not transfer between the cases. For a stochastic
regime model, implement fixed-path HMC plus Metropolis or PGAS path moves
(Section~\ref{sec:pgas}, Section~\ref{sec:joint-hmc}) and validate the invariant distribution
against exact enumeration. For a deterministic threshold model, implement a
non-gradient joint-state or particle-MCMC baseline and verify that regime-only
moves are rejected exactly as the support calculation of Section~\ref{sec:targets}
predicts; a sampler that accepts them is sampling a different target. For a
continuous kink target, the pseudo-marginal and joint-HMC routes of
Section~\ref{sec:pmcmc} are both exact, and their agreement on matched synthetic data
sets is the natural acceptance criterion; solver-branch value-continuity
tests per Section~\ref{sec:solver} are a precondition for trusting any gradient route.

\subsection*{Event-aware HMC only for a genuine switching surface}

Reflection/refraction and discontinuous HMC (Section~\ref{sec:event}, Section~\ref{sec:dhmc})
earn their complexity only when the first stage has identified an explicit
continuous switching surface with a genuine density jump across it. In that
case, test first-crossing accuracy, one-sided full-posterior evaluation,
energy conservation at reflections and refractions, reversibility, and
behavior under boundary grazing. If the boundary locator fails on the actual
model, the event mechanism has no boundary to work with, and smoothing the
surface to rescue it substitutes a different target, as
Section~\ref{sec:relaxations} makes precise; the approach itself has to be
abandoned or the boundary representation repaired.

\subsection*{Test value continuity before trusting gradients}

Finally, at every scale, the evaluation-level tests of Section~\ref{sec:solver},
the branch-boundary bound \eqref{eq:98}, branch-stable derivatives, and
deterministic tie-breaking, are prerequisites for interpreting any
gradient-based sampler's diagnostics. Without them, a rejection cluster is
uninterpretable: it may be model geometry, or it may be a solver branch the
sampler was never told about.

\section{Conclusions and limits}
\label{sec:conclusion}

The literature gives useful exact building blocks, not a finished
nonlinear-ZLB HMC methodology. OccBin and the piecewise Kalman filter solve
the conditional linear problem. Aruoba et al. show how a piecewise-linear
regime mixture can support an efficient particle likelihood. Afshar--Domke and
Nishimura--Dunson--Lu show how Hamiltonian motion can cross true energy
discontinuities; Zhou shows how to preserve a mixed discrete-continuous target
without an artificial ordinal embedding. Childers et al. show the efficiency
of joint latent-state HMC for smooth differentiable state-space models, while
Alenl\"ov et al. state the additional conditions needed for pseudo-marginal
HMC.

The defensible direction is therefore conditional on the model's support. A
deterministic ZLB calls for a linear piecewise-Kalman benchmark, a nonlinear
particle authority, and eventually an event-aware continuous kernel. A
stochastic regime calls for a mixed target, fixed-path HMC, and PGAS or mixed
HMC. Softplus and Concrete relaxations are useful comparison arms but change
the target.

The two case studies sharpen this conditional into two concrete programs. For
the shadow-rate term-structure model of Section~\ref{sec:shadow}, the hard-bound
observation-map construction defines a ladder of targets whose hard-bound
variants are continuous kink targets: the exact routes are the bootstrap
hard-bound particle filter with PMMH and joint kinked-target HMC, which must
agree on matched synthetic data; the fully adapted cell decomposition is a
proposal and diagnostic layer conditional on controlled polytope-probability
and truncated-draw numerics; the softplus and UKF routes remain declared
approximations with stated bias bounds and a sensitivity protocol; and the
closure-filter shadow-rate line from Black through Krippner, Wu--Xia,
Kim--Priebsch, Christensen--Rudebusch, and Opschoor--van der Wel is the domain
benchmark family, with Pericoli and Taboga's closure-free data-augmentation
MCMC the standard against which any likelihood-fidelity claim must be
measured. For the New Keynesian model of Section~\ref{sec:nk}, the worked
linear-complementarity analysis shows the genuine discontinuity lives in the
solution correspondence, nonexistence and pervasive multiplicity of
verified paths, so the first modeling act must be a declared selection law,
deterministic or sunspot-stochastic, and the mixed-support kernels apply only
under the stochastic completion.

A final limitation: no literature survey
establishes the posterior correctness of any particular implementation; that
requires the validation program sketched in Section~\ref{sec:guidance}: a defined
target, a tied-out linear benchmark, a particle authority checked against
enumeration, exact kernels validated against that authority, and
value-continuity tests on the evaluation code, all carried out on the model
actually being estimated.


\section*{References}
\addcontentsline{toc}{section}{References}
\begin{list}{}{\setlength{\leftmargin}{2em}\setlength{\itemindent}{-2em}\setlength{\itemsep}{2pt}\setlength{\parsep}{0pt}}

\item Afshar, Hadi Mohasel, and Justin Domke. 2015. ``Reflection, Refraction, and Hamiltonian Monte Carlo.'' \emph{Advances in Neural Information Processing Systems} 28: 3007--3015. \url{http://hdl.handle.net/1885/103828}.

\item Alenl\"ov, Johan, Arnaud Doucet, and Fredrik Lindsten. 2021. ``Pseudo-Marginal Hamiltonian Monte Carlo.'' \emph{Journal of Machine Learning Research} 22(141): 1--45. \url{https://jmlr.org/papers/v22/19-486.html}.

\item Allik, Bethany, Cory Miller, Michael J. Piovoso, and Ryan Zurakowski. 2016. ``The Tobit Kalman Filter: An Estimator for Censored Measurements.'' \emph{IEEE Transactions on Control Systems Technology} 24(1): 365--371. \url{https://doi.org/10.1109/TCST.2015.2432155}.

\item Alspach, Daniel L., and Harold W. Sorenson. 1972. ``Nonlinear Bayesian Estimation Using Gaussian Sum Approximations.'' \emph{IEEE Transactions on Automatic Control} 17(4): 439--448. \url{https://doi.org/10.1109/TAC.1972.1100034}.

\item Andrieu, Christophe, Arnaud Doucet, and Roman Holenstein. 2010. ``Particle Markov Chain Monte Carlo Methods.'' \emph{Journal of the Royal Statistical Society: Series B} 72(3): 269--342. \url{https://doi.org/10.1111/j.1467-9868.2009.00736.x}.

\item Aruoba, S. Bora\u{g}an, Pablo Cuba-Borda, Kenji Higa-Flores, Frank Schorfheide, and Sergio Villalvazo. 2021. ``Piecewise-Linear Approximations and Filtering for DSGE Models with Occasionally Binding Constraints.'' Federal Reserve Bank of Philadelphia Working Paper 20-13. \url{https://doi.org/10.21799/frbp.wp.2020.13}.

\item Aruoba, S. Bora\u{g}an, Pablo Cuba-Borda, and Frank Schorfheide. 2018. ``Macroeconomic Dynamics Near the ZLB: A Tale of Two Countries.'' \emph{Review of Economic Studies} 85(1): 87--118. \url{https://doi.org/10.1093/restud/rdx027}.

\item Bauer, Michael D., and Glenn D. Rudebusch. 2016. ``Monetary Policy Expectations at the Zero Lower Bound.'' \emph{Journal of Money, Credit and Banking} 48(7): 1439--1465. \url{https://doi.org/10.1111/jmcb.12338}.

\item Black, Fischer. 1995. ``Interest Rates as Options.'' \emph{Journal of Finance} 50(5): 1371--1376. \url{https://doi.org/10.1111/j.1540-6261.1995.tb05182.x}.

\item Blom, Henk A. P., and Yaakov Bar-Shalom. 1988. ``The Interacting Multiple Model Algorithm for Systems with Markovian Switching Coefficients.'' \emph{IEEE Transactions on Automatic Control} 33(8): 780--783. \url{https://doi.org/10.1109/9.1299}.

\item Boehl, Gregor. 2022. ``Efficient Solution and Computation of Models with Occasionally Binding Constraints.'' \emph{Journal of Economic Dynamics and Control} 143: 104523. \url{https://doi.org/10.1016/j.jedc.2022.104523}.

\item Boehl, Gregor, and Felix Strobel. 2023. ``Estimation of DSGE Models with the Effective Lower Bound.'' \emph{Journal of Economic Dynamics and Control} 158: 104784. \url{https://doi.org/10.1016/j.jedc.2023.104784}.

\item Botev, Zdravko I. 2017. ``The Normal Law Under Linear Restrictions: Simulation and Estimation via Minimax Tilting.'' \emph{Journal of the Royal Statistical Society: Series B} 79(1): 125--148. \url{https://doi.org/10.1111/rssb.12162}. [Cited from publication metadata.]

\item Cha\^ari, Lotfi, Jean-Yves Tourneret, Caroline Chaux, and Hadj Batatia. 2016. ``A Hamiltonian Monte Carlo Method for Non-Smooth Energy Sampling.'' \emph{IEEE Transactions on Signal Processing} 64(21): 5585--5594. \url{https://doi.org/10.1109/TSP.2016.2585120}.

\item Childers, David, Jes\'us Fern\'andez-Villaverde, Jesse Perla, Christopher Rackauckas, and Peifan Wu. 2022. ``Differentiable State-Space Models and Hamiltonian Monte Carlo Estimation.'' NBER Working Paper 30573. \url{https://doi.org/10.3386/w30573}.

\item Christensen, Jens H. E., and Glenn D. Rudebusch. 2015. ``Estimating Shadow-Rate Term Structure Models with Near-Zero Yields.'' \emph{Journal of Financial Econometrics} 13(2): 226--259. \url{https://doi.org/10.1093/jjfinec/nbu010}.

\item Cuba-Borda, Pablo, Luca Guerrieri, Matteo Iacoviello, and Molin Zhong. 2019. ``Likelihood Evaluation of Models with Occasionally Binding Constraints.'' \emph{Journal of Applied Econometrics} 34(7): 1073--1085. \url{https://doi.org/10.1002/jae.2733}. [Cited from publication metadata.]

\item Corenflos, Adrien, James Thornton, George Deligiannidis, and Arnaud Doucet. 2021. ``Differentiable Particle Filtering via Entropy-Regularized Optimal Transport.'' \emph{Proceedings of the 38th International Conference on Machine Learning}, PMLR 139: 2100--2111. \url{https://proceedings.mlr.press/v139/corenflos21a.html}.

\item Deligiannidis, George, Arnaud Doucet, and Michael K. Pitt. 2018. ``The Correlated Pseudomarginal Method.'' \emph{Journal of the Royal Statistical Society: Series B} 80(5): 839--870. \url{https://doi.org/10.1111/rssb.12280}. [Cited from publication metadata.]

\item Garc\'ia-Fern\'andez, \'Angel F., Mark R. Morelande, and Jes\'us Grajal. 2012a. ``Truncated Unscented Kalman Filtering.'' \emph{IEEE Transactions on Signal Processing} 60(7): 3372--3386. \url{https://doi.org/10.1109/TSP.2012.2193393}.

\item Garc\'ia-Fern\'andez, \'Angel F., Mark R. Morelande, and Jes\'us Grajal. 2012b. ``Mixture Truncated Unscented Kalman Filtering.'' In \emph{Proceedings of the 15th International Conference on Information Fusion (FUSION)}, 479--486. Singapore.

\item Geng, Hang, Hongjian Liu, Lifeng Ma, and Xiaojian Yi. 2021. ``Multi-Sensor Filtering Fusion Meets Censored Measurements Under a Constrained Network Environment: Advances, Challenges and Prospects.'' \emph{International Journal of Systems Science} 52(16): 3410--3436. \url{https://doi.org/10.1080/00207721.2021.2005178}.

\item Genz, Alan, and Frank Bretz. 2009. \emph{Computation of Multivariate Normal and t Probabilities}. Lecture Notes in Statistics 195. Berlin: Springer. \url{https://doi.org/10.1007/978-3-642-01689-9}. [Cited from publication metadata.]

\item Geweke, John. 2004. ``Getting It Right: Joint Distribution Tests of Posterior Simulators.'' \emph{Journal of the American Statistical Association} 99(467): 799--804. \url{https://doi.org/10.1198/016214504000001132}. [Cited from publication metadata.]

\item Giovannini, Massimo, Philipp Pfeiffer, and Marco Ratto. 2021. ``Efficient and Robust Inference of Models with Occasionally Binding Constraints.'' JRC Working Papers in Economics and Finance 2021/3. \url{https://hdl.handle.net/10419/249365}.

\item Gokce, Murat, and Mustafa Kuzuoglu. 2015. ``Unscented Kalman Filter-Aided Gaussian Sum Filter.'' \emph{IET Radar, Sonar \& Navigation} 9(5): 589--599. \url{https://doi.org/10.1049/iet-rsn.2014.0088}.

\item Gordon, Neil J., David J. Salmond, and Adrian F. M. Smith. 1993. ``Novel Approach to Nonlinear/Non-Gaussian Bayesian State Estimation.'' \emph{IEE Proceedings F (Radar and Signal Processing)} 140(2): 107--113. \url{https://doi.org/10.1049/ip-f-2.1993.0015}. [Cited from publication metadata.]

\item Gorovoi, Viatcheslav, and Vadim Linetsky. 2004. ``Black's Model of Interest Rates as Options, Eigenfunction Expansions and Japanese Interest Rates.'' \emph{Mathematical Finance} 14(1): 49--78. \url{https://doi.org/10.1111/j.0960-1627.2004.00181.x}. [Cited from publication metadata.]

\item Guerrieri, Luca, and Matteo Iacoviello. 2015. ``OccBin: A Toolkit for Solving Dynamic Models with Occasionally Binding Constraints Easily.'' \emph{Journal of Monetary Economics} 70: 22--38. \url{https://doi.org/10.1016/j.jmoneco.2014.08.005}.

\item Holden, Tom D. 2016. ``Computation of Solutions to Dynamic Models with Occasionally Binding Constraints.'' Unpublished working paper, University of Surrey. \url{https://github.com/tholden/dynareOBC}.

\item Holden, Tom D., and Michael Paetz. 2012. ``Efficient Simulation of DSGE Models with Inequality Constraints.'' Quantitative Macroeconomics Working Papers 21207b, Hamburg University. [Cited from publication metadata; the original anticipated-shock approach.]

\item Holden, Tom D. 2017. ``Tractable Estimation and Smoothing of Highly Nonlinear Dynamic State-Space Models.'' Unpublished working paper, University of Surrey.

\item Holden, Tom D. 2023. ``Existence and Uniqueness of Solutions to Dynamic Models with Occasionally Binding Constraints.'' \emph{Review of Economics and Statistics} 105(6): 1481--1499. \url{https://doi.org/10.1162/rest_a_01122}.

\item Kandepu, Rambabu, Lars Imsland, and Bjarne A. Foss. 2008. ``Constrained State Estimation Using the Unscented Kalman Filter.'' In \emph{Proceedings of the 16th Mediterranean Conference on Control and Automation}, 1453--1458. Ajaccio: IEEE. \url{https://doi.org/10.1109/MED.2008.4602001}.

\item Kim, Don H., and Marcel Priebsch. 2013. ``Estimation of Multi-Factor Shadow-Rate Term Structure Models.'' Preliminary draft, October 9. Washington: Board of Governors of the Federal Reserve System.

\item Krippner, Leo. 2012. ``Modifying Gaussian Term Structure Models When Interest Rates Are Near the Zero Lower Bound.'' Reserve Bank of New Zealand Discussion Paper DP2012/02.

\item Lemke, Wolfgang, and Andreea Liliana Vladu. 2017. ``Below the Zero Lower Bound: A Shadow-Rate Term Structure Model for the Euro Area.'' ECB Working Paper No. 1991. Frankfurt: European Central Bank.

\item Lindsten, Fredrik, Michael I. Jordan, and Thomas B. Sch\"on. 2014. ``Particle Gibbs with Ancestor Sampling.'' \emph{Journal of Machine Learning Research} 15: 2145--2184. \url{https://jmlr.org/papers/v15/lindsten14a.html}.

\item Lubik, Thomas A., and Frank Schorfheide. 2003. ``Computing Sunspot Equilibria in Linear Rational Expectations Models.'' \emph{Journal of Economic Dynamics and Control} 28(2): 273--285. \url{https://doi.org/10.1016/S0165-1889(02)00153-7}.

\item Neal, Radford M. 2011. ``MCMC Using Hamiltonian Dynamics.'' In \emph{Handbook of Markov Chain Monte Carlo}, edited by Steve Brooks, Andrew Gelman, Galin L. Jones, and Xiao-Li Meng, 113--162. Chapman and Hall/CRC. \url{https://doi.org/10.1201/b10905-6}.

\item Nishimura, Akihiko, David B. Dunson, and Jianfeng Lu. 2020. ``Discontinuous Hamiltonian Monte Carlo for Discrete Parameters and Discontinuous Likelihoods.'' \emph{Biometrika} 107(2): 365--380. \url{https://doi.org/10.1093/biomet/asz083}.

\item Opschoor, Daan, and Michel van der Wel. 2024. ``A Smooth Shadow-Rate Dynamic Nelson-Siegel Model for Yields at the Zero Lower Bound.'' \emph{Journal of Business \& Economic Statistics} 43(2): 298--311. \url{https://doi.org/10.1080/07350015.2024.2365779}.

\item Pakman, Ari, and Liam Paninski. 2014. ``Exact Hamiltonian Monte Carlo for Truncated Multivariate Gaussians.'' \emph{Journal of Computational and Graphical Statistics} 23(2): 518--542. \url{https://doi.org/10.1080/10618600.2013.788448}.

\item Pericoli, Marcello, and Marco Taboga. 2018. ``Nearly Exact Bayesian Estimation of Non-Linear No-Arbitrage Term-Structure Models.'' Banca d'Italia Temi di discussione (Working Papers) 1189, September 2018. [We cite the 2018 working-paper version.]

\item Poyiadjis, George, Arnaud Doucet, and Sumeetpal S. Singh. 2011. ``Particle Approximations of the Score and Observed Information Matrix in State Space Models with Application to Parameter Estimation.'' \emph{Biometrika} 98(1): 65--80. \url{https://doi.org/10.1093/biomet/asq062}.

\item Priebsch, Marcel A. 2013. ``Computing Arbitrage-Free Yields in Multi-Factor Gaussian Shadow-Rate Term Structure Models.'' Finance and Economics Discussion Series 2013-63. Washington: Board of Governors of the Federal Reserve System. \url{https://doi.org/10.17016/feds.2013.63}.

\item \'Scibior, Adam, and Frank Wood. 2021. ``Differentiable Particle Filtering without Modifying the Forward Pass.'' arXiv:2106.10314v2. \url{https://arxiv.org/abs/2106.10314}.

\item Talts, Sean, Michael Betancourt, Daniel Simpson, Aki Vehtari, and Andrew Gelman. 2018. ``Validating Bayesian Inference Algorithms with Simulation-Based Calibration.'' arXiv:1804.06788. \url{https://arxiv.org/abs/1804.06788}. [Cited from publication metadata.]

\item Teixeira, Bruno O. S., Leonardo A. B. T\^orres, Luis A. Aguirre, and Dennis S. Bernstein. 2010. ``On Unscented Kalman Filtering with State Interval Constraints.'' \emph{Journal of Process Control} 20(1): 45--57. \url{https://doi.org/10.1016/j.jprocont.2009.10.007}.

\item Torgander, Jakob, M\aa{}ns Magnusson, and Jonas Wallin. 2024. ``Hamiltonian Monte Carlo with Categorical Parameters Using the Concrete Distribution.'' Workshop at the 6th Symposium on Advances in Approximate Bayesian Inference, non-archival, 1--12.

\item van der Merwe, Rudolph, Arnaud Doucet, Nando de Freitas, and Eric A. Wan. 2000. ``The Unscented Particle Filter.'' In \emph{Advances in Neural Information Processing Systems} 13: 584--590.

\item Wang, Yanbo, Fasheng Wang, Jianjun He, and Fuming Sun. 2020. ``Iterative Truncated Unscented Particle Filter.'' \emph{Information} 11(4): 214. \url{https://doi.org/10.3390/info11040214}.

\item Wu, Jing Cynthia, and Fan Dora Xia. 2016. ``Measuring the Macroeconomic Impact of Monetary Policy at the Zero Lower Bound.'' \emph{Journal of Money, Credit and Banking} 48(2--3): 253--291. \url{https://doi.org/10.1111/jmcb.12300}.

\item Zhang, Hongwei, Xiaohu Zhang, Weixin Xie, and Xia Yang. 2020. ``Constrained Multiple Model UK Filter.'' In \emph{Proceedings of the 15th IEEE International Conference on Signal Processing (ICSP)}, 48--51. Beijing: IEEE. \url{https://doi.org/10.1109/ICSP48669.2020.9320991}.

\item Zhou, Guangyao. 2020. ``Mixed Hamiltonian Monte Carlo for Mixed Discrete and Continuous Variables.'' \emph{Advances in Neural Information Processing Systems} 33. \url{https://arxiv.org/abs/1909.04852}.

\end{list}


\end{document}
